% !TEX program = xelatex
% !TEX encoding = UTF-8 Unicode
%
% Introduction to Control Engineering -- lecture notes (English), Fall 2026
% One continuous document; new lectures are folded into the existing sections.
% Everything shared with the Korean edition lives in controlnotes.sty.
% Compile:  xelatex control_notes_EN.tex   (three times, for the table of
%           contents -- publish.sh does that for you)
%
% Rules for what comes next:
%   1. New lecture material is folded into the section of the Part whose topic
%      it belongs to; it is not appended as a new chapter.
%   2. If a lecture opens a genuinely new topic, open a new Part with
%      \partdivider.  Section numbers run on -- they are not reset.
%   3. Put a \lecturebreak where that lecture's new material begins, and add
%      one row to the lecture table in the front matter pointing at the same
%      place.  (lectures.py reads the LAST row of the Korean edition's table to
%      build the manifest's `lectures' field.)
%   4. In "The arc of the course", move the Part concerned to in progress or
%      done; and list a planned item only where there is evidence for it --
%      the plan as the lectures themselves stated it.
%   5. The Part, section and label structure of the two editions is 1:1.
%
\documentclass[11pt]{article}
\usepackage{controlnotes}

%% ---- masthead fields ------------------------------------------------------
%% \lectitle is the old title block's "Lecture notes" plus the range of the
%% lecture/section table below; the term and the date are the other two pieces
%% of the line that used to sit under the author.  \lecnum and \lecturer stay
%% empty: this is one continuous document and no lecturer is named in it.
\coursename{Introduction to Control Engineering}
\term{Fall 2026}
\lectitle{Lecture Notes: Lectures 1--10}
\lecdate{Last updated 10 October 2026}
\scribe{Suhan Woo}

\begin{document}
\makelecheader

\section*{The arc of the course}

In the first lecture the course was laid out in three stages: first
\emph{modelling} dynamical systems --- by transfer functions and in state-variable
form --- and learning to \emph{analyse} the models; then a middle stretch that
leans on linear algebra and on MATLAB; and finally stability, controller design,
and classical control. The destinations named were PID control, state-feedback
controllers, state observers, tracking controllers, and Bode and Nyquist plots.
The Parts of these notes follow that plan.

\begin{itemize}[leftmargin=1.2em,itemsep=1pt,topsep=3pt]
  \item Part I --- Feedback and Modelling \hfill {\small\S\S1--2}
  \item Part II --- State-Space Models and Linearization \hfill {\small\S\S3--4}
  \item Part III --- Transfer Functions \hfill {\small\S5}
  \item Part IV --- Frequency Response and Stability \hfill {\small\S\S6--7}
  \item Part V --- Linear Algebra for State Space \hfill {\small\S\S8--11, in progress}
  \item \emph{Stability, controller design, and classical control} --- state-feedback
        controllers, state observers, tracking controllers, PID control
        \hfill {\small planned}
\end{itemize}

The planned items are listed as the first lecture described them; how they divide
into Parts will be settled as the lectures arrive. The goal of
Question~\ref{q:goal} --- bringing the two carts to rest $1$\,m ahead after
$5$\,s without knowing their parameters --- is a controller-design problem; the
Parts so far build the models and the analysis that design will need.

\section*{How these notes are organised}

These notes are ordered by topic rather than by date. They are divided into Parts,
each a stage of the course, and every new lecture is folded into the Part and
section it belongs to, so nothing is said twice. Section numbers run on through
the whole document; they do not restart with each Part. A hairline rule carrying a
lecture number and date marks the point where each lecture's new material begins,
and the table below lists the same points; where several lectures are written up
together as one unit, a single rule carries the whole range. \textbf{Remarks} and worked
\textbf{Examples} fill in derivations, examples and figures beyond what fitted on
the board, and sit with the topic they belong to --- so a rule marks where a
lecture began, not a claim that everything after it was said in that lecture.

\smallskip
\begin{center}\small
\begin{tabular}{@{}lll@{}}
\toprule
Lecture & Date & Begins at\\
\midrule
1 & Wed, 2 Sep 2026  & \S1\\
2 & Mon, 7 Sep 2026  & \S2\\
3 & Wed, 9 Sep 2026  & \S4\\
4 & Mon, 14 Sep 2026 & \S5\\
5 & Wed, 16 Sep 2026 & \S\ref{sec:polezero}\\
6 & Mon, 21 Sep 2026 & \S\ref{sec:second}\\
7--10 & Wed, 23 Sep -- Wed, 7 Oct 2026 & \S\ref{sec:vspace}\\
\bottomrule
\end{tabular}
\end{center}

{\small\tableofcontents}

%% ================================================================ Part I
\partdivider{Part I --- Feedback and Modelling}

\lecturebreak{Lecture 1}{2 September 2026}

%%%%%%%%%%%%%%%%%%%%%%%%%%%%%%%%%%%%%%%%%%%%%%%%%%%%%%%%%%%%%%%%%%%%%%%%%%%%%%
\section{What Control Engineering Is About}

\subsection{The question} Control engineering asks one question:
\begin{center}
\emph{how should the input be chosen so that the output does what we want?}
\end{center}
To answer it we must first know how the input and the output are related, and for
a physical system that relation is almost always a \emph{differential equation}.
So the course begins with differential equations, not with controllers.

Two words will follow us all semester. The first is \textbf{uncertainty}: the
differential equation we write down is never exactly the one the world obeys. The
second is \textbf{feedback}: measuring the output and feeding it back into the
input is what makes a design survive that uncertainty.

A control system in practice is assembled from sensors, actuators, computers, the
controller itself, and the communication between them. In these notes all of that
is idealised away into a signal $u(t)$ going in and a signal $y(t)$ coming out.
The \emph{output} $y$ need not be a continuously varying quantity to be tracked;
it may equally well be ``get to that point and stay there''. Both are handled by
the same machinery.

\subsection{Dynamic versus static} A \emph{static} relation determines $y$ at time
$t$ from $u$ at the same time $t$. A \emph{dynamic} relation has memory: the
present output depends on the past input. The simplest example is a first-order
linear differential equation,
\begin{equation}\label{eq:rl}
   L\dot y + Ry = u(t),
\end{equation}
which is, for instance, an $RL$ circuit driven by a voltage source $u$, with $y$
the current. We write $y^{*}(t)$ for the \emph{desired} output --- the thing the
controller is supposed to produce --- and $u(t)$ for the control input.

\begin{note}
The symbol $y$ is used for whatever the output happens to be: a current in
\eqref{eq:rl}, a position in the next example. Keep track of which is which; the
mathematics does not care, but the units do.
\end{note}

\subsection{A first example: holding a mass up against gravity}
Consider a mass $m$ that can move vertically, with position $y(t)$ measured
upwards, acted on by a control force $u(t)$ and by gravity
(Figure~\ref{fig:mass}):
\begin{equation}\label{eq:mass}
   m\ddot y = u - mg .
\end{equation}
We want to hold it at height $h$: the desired output is the constant
$y^{*}(t)=h$, and the mass starts there at rest,
$y(0)=h$, $\dot y(0)=0$.

\begin{figure}[htb]
\centering
\begin{tikzpicture}[>=Stealth,scale=1.0]
  \draw[thick] (-0.15,0) -- (5.0,0);
  \foreach \x in {0.05,0.35,...,4.95}{\draw[thin] (\x,0) -- ++(-0.22,-0.22);}
  \draw[->,thick] (0.25,0) -- (0.25,3.55) node[left] {$y$};
  \draw[dashed] (0.25,2.50) -- (4.95,2.50);
  \node[right] at (4.95,2.50) {$y^{*}=h$};
  \draw[thick,fill=black!5] (1.95,2.10) rectangle (3.15,2.90);
  \node at (2.55,2.50) {$m$};
  \draw[very thick,->] (2.20,1.15) -- (2.20,2.07);
  \node[left,sig] at (2.10,1.58) {$u(t)$};
  \draw[very thick,->] (2.90,2.07) -- (2.90,1.15);
  \node[right,sig] at (2.92,1.58) {$mg$};
  \draw[<->] (1.15,0) -- (1.15,2.50);
  \node[left,sig] at (1.13,1.25) {$y(t)$};
\end{tikzpicture}
\caption{A mass held against gravity by a control force. The control problem is
to choose $u(t)$ so that $y(t)$ stays at $h$.}
\label{fig:mass}
\end{figure}

\smallskip
\noindent\textbf{(a) Open loop.} Setting $\ddot y=0$ in \eqref{eq:mass} gives
$u=mg$. With that constant input $\ddot y \equiv 0$, and with the initial
conditions above $y(t)\equiv h$ for all time. On paper, the problem is solved
without any feedback at all.

\smallskip
\noindent\textbf{(b) The trouble: uncertainty.} We do not know $m$. We know an
estimate $\hat m$, and so we apply $u=\hat m g$. Now \eqref{eq:mass} reads
$m\ddot y = \hat m g - mg = (\hat m - m)g$, a \emph{constant nonzero
acceleration}. Integrating twice,
\[
   y(t) = h + \frac{(\hat m-m)g}{2m}\,t^{2}\;\longrightarrow\;\pm\infty .
\]
However small the modelling error $\hat m - m$, the mass eventually leaves. The
open-loop design is not merely inaccurate; it fails qualitatively.

\smallskip
\noindent\textbf{(c) The cure: feedback.} Measure $y$ and $\dot y$ and use
\begin{equation}\label{eq:pd}
   u(t) = -c\,\dot y - k\,(y-h) + \hat m g ,
\end{equation}
where $c$ and $k$ are \emph{gains}: numbers we choose. They carry no uncertainty,
because we are the ones who pick them. Substituting \eqref{eq:pd} into
\eqref{eq:mass},
\begin{equation}\label{eq:massfb}
   m\ddot y + c\,\dot y + k\,(y-h) = (\hat m-m)g .
\end{equation}
For $c,k>0$ this is a damped second-order system driven by a constant, so $y$
settles, and setting $\ddot y=\dot y=0$ gives the resting value
\begin{equation}\label{eq:offset}
   y(\infty) = h + \frac{(\hat m - m)g}{k}.
\end{equation}
The error is now \emph{bounded}, and it shrinks as $k$ grows. Feedback did not
remove the uncertainty; it converted an unbounded error into a small one.

\begin{question}
Why does this work? What are good values of $c$ and $k$? How large may $k$ be
before something else goes wrong? Answering these is what the rest of the course
is for.
\end{question}

\subsection{Uncertainty, robustness, and learning}
Designing a controller for the nominal model is only half the job; the other half
is guaranteeing that it still behaves over the whole \emph{set} of plants that the
real system might be. That is the subject of \emph{robust} control engineering.

As an aside, the lecture contrasted this with reinforcement learning, where a
policy is trained from a reward signal rather than from a model. Its
characteristic difficulty is the \emph{sim-to-real gap}: a policy tuned in
simulation need not work on hardware. The standard remedies --- domain
randomisation and friends, i.e.\ training over a randomised family of simulators
so that no single mismatch is fatal --- are robustness ideas in disguise. They are
heuristic, but they are valuable exactly where a usable physical model is hard to
come by. A model-based design, by contrast, is mathematical: one can in principle
reason backwards from a specification to a controller. Even that is not easy, and
much of this course is about learning how.

%% ---------------------------------------------------------------- Lecture 2
\lecturebreak{Lecture 2}{7 September 2026}

%%%%%%%%%%%%%%%%%%%%%%%%%%%%%%%%%%%%%%%%%%%%%%%%%%%%%%%%%%%%%%%%%%%%%%%%%%%%%%
\section{Modelling: the Two-Cart System}

\subsection{The running example}
The example that will accompany us through the course is the \emph{two-cart
system} of Figure~\ref{fig:twocart}: a cart of mass $m_1$ at position $p_1(t)$,
pushed by a force $u(t)$, coupled through a spring of stiffness $k$ and a damper
$b$ to a second cart of mass $m_2$ at position $p_2(t)$. Both carts feel ground
friction with coefficient $c$. The measured output is the position of the second
cart, $y(t) = p_2(t)$. We take $p_1(0)=p_2(0)=0$: positions are measured from
where the carts start.

\begin{figure}[htb]
\centering
\begin{tikzpicture}[>=Stealth,scale=1.05]
  \draw[thick] (-0.85,0) -- (7.75,0);
  \foreach \x in {-0.6,-0.25,...,7.7}{\draw[thin] (\x,0) -- ++(-0.25,-0.25);}
  \draw[thick,fill=black!5] (0.40,0.36) rectangle (2.00,1.60);
  \node at (1.20,0.98) {$m_1$};
  \draw[thick,fill=black!5] (5.00,0.36) rectangle (6.60,1.60);
  \node at (5.80,0.98) {$m_2$};
  \foreach \x in {0.80,1.60,5.40,6.20}{\draw[thick,fill=white] (\x,0.18) circle (0.18);}
  \draw[thick] (2.00,1.28) -- (2.35,1.28);
  \draw[thick,decorate,decoration={coil,aspect=0.62,segment length=3.8mm,amplitude=2.2mm}]
       (2.35,1.28) -- (4.65,1.28);
  \draw[thick] (4.65,1.28) -- (5.00,1.28);
  \node[above] at (3.50,1.44) {$k$};
  \draw[thick] (2.00,0.68) -- (3.20,0.68);
  \draw[thick] (3.20,0.46) rectangle (3.78,0.90);
  \draw[thick] (3.62,0.46) -- (3.62,0.90);
  \draw[thick] (3.62,0.68) -- (5.00,0.68);
  \node[below] at (2.58,0.63) {$b$};
  \draw[very thick,->] (-0.80,0.98) -- (0.38,0.98);
  \node[above] at (-0.22,1.02) {$u(t)$};
  \draw[dashed,gray] (1.20,0) -- (1.20,-0.62);
  \draw[->] (1.20,-0.62) -- (0.35,-0.62);
  \node[right,lbl] at (1.26,-0.62) {$c\,\dot p_1$};
  \draw[dashed,gray] (5.80,0) -- (5.80,-0.62);
  \draw[->] (5.80,-0.62) -- (4.95,-0.62);
  \node[right,lbl] at (5.86,-0.62) {$c\,\dot p_2$};
  \draw[dashed] (1.20,1.60) -- (1.20,2.42);
  \draw[dashed] (5.80,1.60) -- (5.80,2.42);
  \draw[->] (1.20,2.22) -- ++(0.90,0); \node[above] at (1.65,2.22) {$p_1(t)$};
  \draw[->] (5.80,2.22) -- ++(0.90,0); \node[above] at (6.25,2.22) {$p_2(t)=y(t)$};
\end{tikzpicture}
\caption{The two-cart system. The spring $k$ and the damper $b$ act between the
carts; the ground friction $c$ acts on each cart separately and opposes its
velocity.}
\label{fig:twocart}
\end{figure}

\begin{question}[the goal for the semester]\label{q:goal}
Make the system move forward and \emph{come to rest $1$\,m ahead after $5$
seconds} --- and do it without knowing the values of $m_1$, $k$, $m_2$.
\end{question}

The second half of Question~\ref{q:goal} is the real content. Anyone can hit a
target when the plant is known exactly; a controller that works without that
knowledge is a control design.

\subsection{Equations of motion}
Newton's law for each cart. Cart~1 is pushed by $u$, pulled by the spring towards
cart~2, damped by the relative motion of the two carts, and dragged by ground
friction; cart~2 sees the same, minus the input:
\begin{equation}\label{eq:eom}
\left\{
\begin{aligned}
   m_1\ddot p_1 &= u(t) + k\,(p_2-p_1) + b\,(\dot p_2-\dot p_1) - c\,\dot p_1,\\
   m_2\ddot p_2 &= \phantom{u(t)\;+{}} k\,(p_1-p_2) + b\,(\dot p_1-\dot p_2) - c\,\dot p_2,\\
   y &= p_2 .
\end{aligned}\right.
\end{equation}

\begin{remark}[the sign of the coupling terms]
A damper mounted in parallel with the spring resists the \emph{relative} motion of
its two ends, so its force on cart~1 has exactly the same sign pattern as the
spring force: $+\,k(p_2-p_1)$ and $+\,b(\dot p_2-\dot p_1)$, and the mirror images
on cart~2. A quick way to see that these signs are the physical ones is to look at
the stretch of the spring, $q=p_1-p_2$. Subtracting the two equations (with $c=0$
for simplicity),
\[
   \ddot q + b\Bigl(\tfrac1{m_1}+\tfrac1{m_2}\Bigr)\dot q
           + k\Bigl(\tfrac1{m_1}+\tfrac1{m_2}\Bigr)q = \frac{u}{m_1},
\]
a damped oscillator: the internal mode, in which the carts swing against each
other, loses energy through the damper, as it should. Flipping the sign of either
coupling term would give that mode negative stiffness or negative damping.
\end{remark}

\begin{note}[symbols]
Throughout, $c$ is the \emph{ground friction} coefficient and $b$ is the
\emph{damper between the carts}.
\end{note}

\subsection{Three ways to describe the same system}
\begin{figure}[htb]
\centering
\begin{tikzpicture}[>=Stealth]
  \node (ph) [blk] {physical\\ system};
  \node (de) [blk,right=17mm of ph] {differential\\ equation};
  \node (tf) [blk,above right=6mm and 20mm of de] {transfer function\\ $P(s)=Y(s)/U(s)$};
  \node (ss) [blk,below right=6mm and 20mm of de] {state space\\ $\dot x=Ax+Bu$};
  \draw[->,thick] (ph) -- node[above,lbl] {modelling} (de);
  \draw[->,thick] (de.east) -- node[above,sloped,lbl,pos=0.60] {Laplace} (tf.west);
  \draw[->,thick] (de.east) -- node[below,sloped,lbl,pos=0.60] {state variables} (ss.west);
  \draw[<->,thick] (tf) -- node[right,lbl,align=left]
        {\S\ref{sec:ss2tf}\\ realization} (ss);
  \draw[->,dashed] (ph.north) to[bend left=26]
        node[above,lbl] {data-driven} (tf.north west);
  \draw[->,dashed] (ph.south) to[bend right=26]
        node[below,lbl] {data-driven} (ss.south west);
\end{tikzpicture}
\caption{Routes from a physical system to a mathematical model. The solid path is
this course. The dashed ones are system identification / learning from data
(neural networks and the like), which skip the differential equation altogether.}
\label{fig:routes}
\end{figure}

The transfer function and the state-space form are two descriptions of the same
differential equation, and (for linear systems) one can pass between them --- the
double arrow in Figure~\ref{fig:routes}. Section~\ref{sec:ss2tf} does one of the
two directions explicitly.

%% =============================================================== Part II
\partdivider{Part II --- State-Space Models and Linearization}

%%%%%%%%%%%%%%%%%%%%%%%%%%%%%%%%%%%%%%%%%%%%%%%%%%%%%%%%%%%%%%%%%%%%%%%%%%%%%%
\section{The State-Space Representation}

\subsection{State variables}
The equations \eqref{eq:eom} are second order. Introduce one variable per
independent piece of ``memory'' the system has --- here, a position and a velocity
per cart:
\[
   x_1 = p_1,\qquad x_2=\dot p_1,\qquad x_3=p_2,\qquad x_4=\dot p_2 .
\]
Then \eqref{eq:eom} becomes four \emph{first-order} equations:
\begin{equation}\label{eq:statelines}
\left\{
\begin{aligned}
   \dot x_1 &= x_2, \\
   \dot x_2 &= \frac{k}{m_1}(x_3-x_1) - \frac{c}{m_1}x_2 + \frac{b}{m_1}(x_4-x_2) + \frac{1}{m_1}u, \\
   \dot x_3 &= x_4, \\
   \dot x_4 &= \frac{k}{m_2}(x_1-x_3) - \frac{c}{m_2}x_4 + \frac{b}{m_2}(x_2-x_4), \\
   y &= x_3 .
\end{aligned}\right.
\end{equation}
The four equations are thoroughly tangled together, but every one of them is
first order, and that is precisely what makes them easy to handle --- by a
computer and by linear algebra alike.

\subsection{The standard form}\label{sec:stdform}
Collecting $x=(x_1,x_2,x_3,x_4)\tp$ we reach the \textbf{standard form of the
state-space representation}:
\begin{equation}\label{eq:std}
   \boxed{\;\dot x(t) = A\,x(t) + B\,u(t),\qquad y(t) = C\,x(t) + D\,u(t).\;}
\end{equation}
Its dimensions are
\begin{equation}\label{eq:dims}
   x(t)\in\R^{n},\qquad u(t)\in\R^{m},\qquad y(t)\in\R^{p},
\end{equation}
so that $A\in\R^{n\times n}$, $B\in\R^{n\times m}$, $C\in\R^{p\times n}$,
$D\in\R^{p\times m}$. The space $\R^{n}$ is called the \textbf{state space}, and
$n$ --- the number of state variables --- is the \emph{order} of the system. A
system with $m=p=1$ is called SISO (single input, single output); otherwise it
is MIMO. Note that the three dimensions are independent: a plant may have many
states and a single input, or several inputs and one output.

For the two-cart system $n=4$, $m=p=1$. Taking $b=0$ to keep the entries short,
\begin{equation}\label{eq:ABCD}
 A=\begin{bmatrix}
     0 & 1 & 0 & 0\\[2pt]
     -\dfrac{k}{m_1} & -\dfrac{c}{m_1} & \dfrac{k}{m_1} & 0\\[6pt]
     0 & 0 & 0 & 1\\[2pt]
     \dfrac{k}{m_2} & 0 & -\dfrac{k}{m_2} & -\dfrac{c}{m_2}
   \end{bmatrix},\quad
 B=\begin{bmatrix}0\\[2pt] \dfrac{1}{m_1}\\[6pt] 0\\[2pt] 0\end{bmatrix},\quad
 C=\begin{bmatrix}0&0&1&0\end{bmatrix},\quad D = 0 .
\end{equation}

\begin{remark}[where the uncertainty lives]
Rows 1 and 3 of $A$ contain no parameters at all. They are not physics; they are
the definitions $\dot x_1=x_2$ and $\dot x_3=x_4$, and they are exactly right
whatever the carts weigh. All the uncertainty of the model is confined to rows 2
and 4. It is a general and useful feature of mechanical state-space models: the
kinematic rows are exact, and only the rows that encode Newton's law carry the
parameters one might not know.
\end{remark}

\begin{remark}[keeping the damper]
If $b$ is not dropped, only rows 2 and 4 change:
\[
   A_{2\bullet}=\Bigl[-\tfrac{k}{m_1},\;-\tfrac{c+b}{m_1},\;\tfrac{k}{m_1},\;\tfrac{b}{m_1}\Bigr],
   \qquad
   A_{4\bullet}=\Bigl[\tfrac{k}{m_2},\;\tfrac{b}{m_2},\;-\tfrac{k}{m_2},\;-\tfrac{c+b}{m_2}\Bigr].
\]
The damper adds to the diagonal velocity terms and couples the two velocities,
exactly as the spring couples the two positions.
\end{remark}

In MATLAB the quadruple is packaged as \verb|sys = ss(A,B,C,D)|.

\subsection{What ``state'' means}
\begin{definition}[state]\label{def:state}
The \emph{state} $x(t)$ is the minimum amount of information about the system at
time $t$ which, together with the input $u(\tau)$ for $\tau\ge t$, determines the
output $y(\tau)$ for all $\tau\ge t$.
\end{definition}
Everything that happened before time $t$ matters only through $x(t)$. For the two
carts, knowing where both carts are and how fast they are going is enough; their
history is irrelevant. The word \emph{minimum} is what makes the definition
useful: it is why the state is four-dimensional here and not five.

\begin{remark}\label{rem:nonlinear}
Definition~\ref{def:state} says nothing about linearity, and indeed the
state-space description works for nonlinear systems just as well
(Section~\ref{sec:nonlinear}). The transfer function does not: it exists only for
linear systems.
\end{remark}

\subsection{Simulating the standard form}
Because \eqref{eq:std} is a system of first-order equations, a computer can march
it forward directly. The crudest such scheme, and the one this course starts
with, is the \emph{forward Euler method}: for a small step $h>0$,
\begin{equation}\label{eq:euler}
   x(t+h) \;\approx\; x(t) + h\,\bigl(Ax(t)+Bu(t)\bigr),
\end{equation}
or, for a general nonlinear right-hand side, $x(t+h)\approx x(t)+h\,f(x(t),u(t))$.
It is nothing more than the first two terms of the Taylor series of $x$ about $t$,
so the error made in one step is of order $h^{2}$. Choosing $h$ is not free,
however: too large a step and the simulation diverges even though the system does
not.

\begin{remark}[how large a step]\label{rem:eulerstab}
The limit on $h$ can be computed exactly on the scalar test equation
$\dot x=\lambda x$, whose true solution decays whenever $\operatorname{Re}\lambda<0$.
One Euler step multiplies $x$ by a fixed factor,
\[
   x_{k+1} = (1+h\lambda)\,x_k ,\qquad x_k = (1+h\lambda)^{k}x_0 ,
\]
so the simulation decays only if $|1+h\lambda|<1$ --- that is, only if $h\lambda$
lies inside the unit disc centred at $-1$. For a real $\lambda<0$ this means
\[
   h < \frac{2}{|\lambda|} .
\]
For the linear system $\dot x=Ax$ the same condition must hold for every
eigenvalue of $A$ at once, so the fastest mode sets the step. An eigenvalue on or
near the imaginary axis --- a lightly damped oscillation --- is the worst case: for
$\lambda=-\sigma\pm j\omega$ with small $\sigma$ the condition $|1+h\lambda|<1$
reads $h<2\sigma/(\sigma^{2}+\omega^{2})$, which is tiny when $\sigma\ll\omega$. A
purely imaginary $\lambda$ can never satisfy it: forward Euler turns every
undamped oscillation into a growing one.
\end{remark}

%% ---------------------------------------------------------------- Lecture 3
\lecturebreak{Lecture 3}{9 September 2026}

%%%%%%%%%%%%%%%%%%%%%%%%%%%%%%%%%%%%%%%%%%%%%%%%%%%%%%%%%%%%%%%%%%%%%%%%%%%%%%
\section{Nonlinear Systems, Equilibria, and Linearization}\label{sec:nonlinear}

\subsection{The nonlinear standard form}
Most real systems are not linear. The general form we will use is
\begin{equation}\label{eq:nl}
   \dot z = f(z,u),\qquad w = h(z,u),
\end{equation}
with $z(t)\in\R^{n}$, $u(t)\in\R^{m}$, $w(t)\in\R^{p}$ and
\[
   f:\R^{n}\times\R^{m}\to\R^{n},\qquad h:\R^{n}\times\R^{m}\to\R^{p}
\]
assumed \emph{regular} --- smooth enough that the derivatives below exist. Written
out, \eqref{eq:nl} is
\[
   \dot z=\begin{pmatrix}\dot z_1\\ \vdots\\ \dot z_n\end{pmatrix}
        =\begin{pmatrix}f_1(z(t),u(t))\\ \vdots\\ f_n(z(t),u(t))\end{pmatrix}.
\]

\begin{note}[notation for this section]
We reserve $(z,u,w)$ for the nonlinear system and $(x,u,y)$ for the linear one,
because the variables of the linearized model will turn out to be
\emph{deviations} from an operating point, not the physical quantities themselves.
Mixing the two up is the single most common error in this material.
\end{note}

\begin{example}[pendulum]\label{ex:pend}
A pendulum of mass $m$ on a massless rod of length $\ell$, with viscous friction
$b$ at the pivot and a torque $u(t)$ applied there (Figure~\ref{fig:pend}), obeys
\[
   m\ell^{2}\ddot\theta = u(t) - mg\ell\sin\theta - b\,\dot\theta,
   \qquad\text{i.e.}\qquad
   m\ell^{2}\ddot\theta + b\,\dot\theta + mg\ell\sin\theta = u(t),
\]
where $\theta$ is measured from the downward vertical. With $z_1=\theta$,
$z_2=\dot\theta$ this is exactly of the form \eqref{eq:nl}:
\begin{equation}\label{eq:pendstate}
\left\{
\begin{aligned}
   \dot z_1 &= z_2,\\
   \dot z_2 &= -\frac{b}{m\ell^{2}}\,z_2 - \frac{g}{\ell}\sin z_1
               + \frac{1}{m\ell^{2}}\,u,\\
   w &= z_1 .
\end{aligned}\right.
\end{equation}
The $\sin z_1$ is the whole difficulty: everything else is linear.
\end{example}

\begin{figure}[htb]
\centering
\begin{tikzpicture}[>=Stealth,scale=1.25]
  \draw[thick] (-1.45,0) -- (1.45,0);
  \foreach \x in {-1.35,-1.05,...,1.42}{\draw[thin] (\x,0) -- ++(-0.22,0.22);}
  \draw[dashed] (0,0) -- (0,-2.95);
  \draw[thick] (0,0) -- (1.30,-2.25);
  \draw[thick,fill=black!10] (1.30,-2.25) circle (0.28);
  \node[sig] at (1.30,-2.25) {$m$};
  \fill (0,0) circle (2.4pt);
  \draw[->] (0,-1.25) arc (-90:-60:1.25);
  \node at (0.38,-1.46) {$\theta$};
  \node[right,sig] at (1.00,-1.66) {$\ell$};
  \draw[very thick,->] (1.30,-2.62) -- (1.30,-3.38);
  \node[right,sig] at (1.34,-3.02) {$mg$};
  \draw[->,thick] (190:0.88) arc (190:256:0.88);
  \node[left,sig] at (-0.94,-0.52) {$u(t)$};
  \node[left,sig] at (-0.06,-0.24) {$b$};
\end{tikzpicture}
\caption{The pendulum of Example~\ref{ex:pend}: $\theta$ is measured from the
downward vertical, $b$ is viscous friction at the pivot, and $u(t)$ is an applied
torque.}
\label{fig:pend}
\end{figure}

\subsection{Equilibria} The trick for dealing with \eqref{eq:nl} is to work near a
point where nothing is happening.

\begin{definition}[equilibrium]\label{def:eq}
A pair $(z^{*},u^{*})$ with $f(z^{*},u^{*}) = 0$ is an \emph{equilibrium}. If the
input is held at the constant value $u^{*}$ and the system starts at
$z(0)=z^{*}$, then $\dot z=0$ and $z(t)\equiv z^{*}$ forever.
\end{definition}
Note that the input must be a \emph{constant} for this to make sense; ``an
equilibrium'' is always an equilibrium \emph{of a particular constant input}.

For the pendulum with $u^{*}=0$: $f=0$ forces $z_2^{*}=0$ and $\sin z_1^{*}=0$,
so $z_1^{*}=0$ (hanging down) or $z_1^{*}=\pi$ (balanced upright). Both are
equilibria; only one of them is one you can leave unattended.

\subsection{The Jacobian}
For $f:\R^{n}\to\R^{n}$ (suppressing $u$ for a moment), the \emph{Jacobian} is the
matrix of all first partial derivatives,
\[
   \Dop f \;=\;
   \begin{pmatrix}
     \dfrac{\partial f_1}{\partial x_1} & \dfrac{\partial f_1}{\partial x_2} & \cdots & \dfrac{\partial f_1}{\partial x_n}\\[10pt]
     \vdots & \vdots & \ddots & \vdots\\[6pt]
     \dfrac{\partial f_n}{\partial x_1} & \dfrac{\partial f_n}{\partial x_2} & \cdots & \dfrac{\partial f_n}{\partial x_n}
   \end{pmatrix}
   \;\in\;\R^{n\times n}.
\]
It is the best linear approximation of $f$ near a point --- the multivariable
analogue of the derivative.

\subsection{Linearization}\label{sec:lin}
Expand $f$ in a Taylor series about the equilibrium $(z^{*},u^{*})$:
\[
   \dot z = f(z,u)
   = \underbrace{f(z^{*},u^{*})}_{=\,0}
   + \underbrace{\left.\frac{\partial f}{\partial z}\right|_{z^{*},u^{*}}}_{=:\,A}
     \underbrace{(z-z^{*})}_{=:\,x}
   + \underbrace{\left.\frac{\partial f}{\partial u}\right|_{z^{*},u^{*}}}_{=:\,B}
     \underbrace{(u-u^{*})}_{=:\,\tilde u}
   + \cdots
\]
The constant term vanishes precisely because we expanded about an equilibrium ---
that is what equilibria are for. Since $z^{*}$ is constant, $\dot x=\dot z$, and
dropping the higher-order terms leaves
\begin{equation}\label{eq:linAB}
   \dot x = A\,x + B\,\tilde u,\qquad
   A = \Dop_{z}f\bigl|_{z^{*},u^{*}},\qquad
   B = \Dop_{u}f\bigl|_{z^{*},u^{*}} .
\end{equation}
The output is treated the same way. From
$w = h(z,u) = h(z^{*},u^{*}) + \Dop_{z}h|_{*}\,x + \Dop_{u}h|_{*}\,\tilde u + \cdots$,
define the output deviation $y := w - h(z^{*},u^{*})$, and then
\begin{equation}\label{eq:linCD}
   y = C\,x + D\,\tilde u,\qquad
   C = \Dop_{z}h\bigl|_{z^{*},u^{*}},\qquad
   D = \Dop_{u}h\bigl|_{z^{*},u^{*}} .
\end{equation}
The result is exactly the standard form \eqref{eq:std}, which is why everything
built for linear systems can be pointed at a nonlinear one --- locally.

\begin{caution}[what the variables mean]
$A,B,C,D$ are \emph{constant} matrices: the Jacobians evaluated at the
equilibrium, not functions of $z$. And $x$, $\tilde u$, $y$ are all
\emph{deviations} --- $x=z-z^{*}$, $\tilde u=u-u^{*}$, $y=w-h(z^{*},u^{*})$.
(It is common to drop the tilde and write $u$ for the input deviation as well.)
If you feed a linearized model the absolute input instead of the deviation,
everything downstream is wrong by a constant.
\end{caution}

\begin{example}[pendulum, continued]
For Example~\ref{ex:pend},
\[
   \Dop_{z}f = \begin{bmatrix} 0 & 1\\[4pt]
     -\dfrac{g}{\ell}\cos z_1 & -\dfrac{b}{m\ell^{2}}\end{bmatrix},
   \qquad
   \Dop_{u}f = \begin{bmatrix}0\\[4pt] \dfrac{1}{m\ell^{2}}\end{bmatrix},
   \qquad
   x = \begin{bmatrix}z_1\\ z_2\end{bmatrix} - z^{*} .
\]
At the downward equilibrium $z^{*}=(0,0)$, $u^{*}=0$, we get $\cos z_1^{*}=1$ and
\[
   A = \begin{bmatrix}0&1\\[3pt] -\dfrac{g}{\ell} & -\dfrac{b}{m\ell^{2}}\end{bmatrix},
   \qquad
   B = \begin{bmatrix}0\\[3pt] \dfrac{1}{m\ell^{2}}\end{bmatrix},
   \qquad C=\begin{bmatrix}1&0\end{bmatrix},\qquad D=0 .
\]
\end{example}

\begin{example}[pendulum, the upright equilibrium]\label{ex:pendup}
The $(2,1)$ entry of $\Dop_z f$ is $-\frac{g}{\ell}\cos z_1^{*}$, so it depends
on where we linearize. At the upright equilibrium $z^{*}=(\pi,0)$, $u^{*}=0$ it
becomes $+g/\ell$:
\[
   A_{\text{up}} = \begin{bmatrix}0&1\\[3pt] \dfrac{g}{\ell} & -\dfrac{b}{m\ell^{2}}\end{bmatrix},
   \qquad
   \det(\lambda I-A_{\text{up}}) = \lambda^{2}+\frac{b}{m\ell^{2}}\lambda-\frac{g}{\ell}.
\]
The constant term of the characteristic polynomial is now negative, so one root is
positive whatever $b$ is --- for $b=0$ the eigenvalues are $\pm\sqrt{g/\ell}$.
The linearization is unstable, as it must be: a pendulum balanced on its end falls
over. At the hanging equilibrium the same polynomial reads
$\lambda^{2}+\frac{b}{m\ell^{2}}\lambda+\frac{g}{\ell}$, whose roots have negative
real part for every $b>0$. One nonlinear model, two equilibria, two qualitatively
different linear models --- which is exactly why an equilibrium must always be
named when a linearization is written down.
\end{example}

\subsection{Where the linear model sits}
The linearized system is not a separate object; it is a window onto the real one.
Figure~\ref{fig:lin} shows the bookkeeping: shift the input by $u^{*}$ on the way
in, and subtract the equilibrium output $h(z^{*},u^{*})$ on the way out.

\begin{figure}[htb]
\centering
\begin{tikzpicture}[>=Stealth]
  \node (sum1) [sumc] {};
  \node (plant) [draw,align=center,inner sep=6pt,minimum height=13mm,
                 right=19mm of sum1] {$\dot z = f(z,u)$\\[2pt] $w = h(z,u)$};
  \node (sum2) [sumc,right=19mm of plant] {};
  \node (in)  [left=14mm of sum1] {$\tilde u$};
  \node (out) [right=14mm of sum2] {$y$};
  \draw[->] (in) -- node[pos=0.62,above,lbl] {$+$} (sum1);
  \draw[->] (sum1) -- node[above,sig] {$u$} (plant);
  \draw[->] (plant) -- node[above,sig] {$w$} (sum2);
  \draw[->] (sum2) -- (out);
  \draw[->] ($(sum1)+(0,13mm)$) node[above,sig] {$u^{*}$}
        -- node[pos=0.72,right,lbl] {$+$} (sum1);
  \draw[->] ($(sum2)+(0,13mm)$) node[above,sig] {$h(z^{*},u^{*})$}
        -- node[pos=0.72,right,lbl] {$-$} (sum2);
  \node[below=2.5mm of plant,lbl] {true nonlinear system};
\end{tikzpicture}
\caption{The linearized model $\dot x=Ax+B\tilde u$, $y=Cx+D\tilde u$
approximates the map $\tilde u\mapsto y$ drawn here, for $\tilde u$ small.}
\label{fig:lin}
\end{figure}

\subsection{Phase portraits}
For a second-order autonomous system
$\dot x_1 = f_1(x_1,x_2)$, $\dot x_2 = f_2(x_1,x_2)$,
one can draw the whole solution family at once: plot the trajectories in the
$(x_1,x_2)$ plane. The picture is called a \emph{phase portrait}. Time itself does
not appear on it, but everything else does --- equilibria are the points where the
vector field vanishes, and by uniqueness of solutions two trajectories never
cross. Figure~\ref{fig:phase} shows a stable spiral.

\begin{figure}[htb]
\centering
\begin{tikzpicture}[>=Stealth,scale=1.2]
  \foreach \a in {-2.4,-1.6,...,2.5}{
    \foreach \bb in {-2.4,-1.6,...,2.5}{
      \pgfmathsetmacro{\dx}{\bb}
      \pgfmathsetmacro{\dy}{-\a-0.7*\bb}
      \pgfmathsetmacro{\nn}{sqrt(\dx*\dx+\dy*\dy)+0.001}
      \draw[->,gray!50,thin] (\a,\bb) -- ++({0.44*\dx/\nn},{0.44*\dy/\nn});
    }}
  \draw[->] (-3.0,0) -- (3.0,0) node[right] {$x_1$};
  \draw[->] (0,-2.95) -- (0,2.95) node[above] {$x_2$};
  \foreach \A in {2.4,1.3}{
    \draw[thick] plot[domain=0:720,samples=100,variable=\t,smooth]
      ({\A*exp(-0.0065216*\t)*(cos(\t)+0.37364*sin(\t))},
       {-1.0676*\A*exp(-0.0065216*\t)*sin(\t)});}
  \foreach \A/\T in {2.4/40, 2.4/180, 1.3/110}{
    \pgfmathsetmacro{\xa}{\A*exp(-0.0065216*\T)*(cos(\T)+0.37364*sin(\T))}
    \pgfmathsetmacro{\ya}{-1.0676*\A*exp(-0.0065216*\T)*sin(\T)}
    \pgfmathsetmacro{\xb}{\A*exp(-0.0065216*(\T+16))*(cos(\T+16)+0.37364*sin(\T+16))}
    \pgfmathsetmacro{\yb}{-1.0676*\A*exp(-0.0065216*(\T+16))*sin(\T+16)}
    \draw[->,thick] (\xa,\ya) -- (\xb,\yb);}
  \foreach \A in {2.4,1.3}{\fill (\A,0) circle (1.7pt);}
  \node[above right,lbl] at (2.42,0.04) {$x(0)$};
  \fill (0,0) circle (2.4pt);
\end{tikzpicture}
\caption{Phase portrait of $\dot x_1=x_2$, $\dot x_2=-x_1-0.7x_2$. Grey arrows are
the vector field; the two heavy curves are trajectories. The black dot at the
origin is the equilibrium, and both trajectories spiral into it.}
\label{fig:phase}
\end{figure}

For a \emph{linear} second-order system $\dot x = Ax$ the whole portrait is fixed
by the two eigenvalues $\lambda_{1,2}$ of $A$, and there are only a handful of
possible pictures:

\begin{center}\small
\renewcommand{\arraystretch}{1.3}
\begin{tabular}{@{}lll@{}}
\toprule
eigenvalues & portrait & behaviour near the equilibrium\\
\midrule
real, both negative          & stable node     & every trajectory enters, without circling\\
real, both positive          & unstable node   & every trajectory leaves\\
real, opposite signs         & saddle          & in along one direction, out along the other\\
complex, $\operatorname{Re}\lambda<0$ & stable focus (spiral) & trajectories spiral in\\
complex, $\operatorname{Re}\lambda>0$ & unstable focus & trajectories spiral out\\
purely imaginary             & centre          & closed orbits, neither in nor out\\
\bottomrule
\end{tabular}
\end{center}

\noindent For real eigenvalues the straight-line trajectories along the
eigenvectors are the skeleton of the picture: the solution is
$x(t)=c_1e^{\lambda_1t}v_1+c_2e^{\lambda_2t}v_2$, and each component grows or
decays along its own eigenvector (Figure~\ref{fig:portraits}). For a complex pair
$\sigma\pm j\omega$ the solution is an $e^{\sigma t}$ envelope around a rotation at
frequency $\omega$ --- exactly the partial-fraction terms of \S\ref{sec:pfe}.

\begin{figure}[htb]
\centering
\begin{tikzpicture}[>=Stealth]
  \foreach \px/\ma/\mb/\mc/\md/\name in {0/-1/0/0/-2.5/stable node,
                                     3.35/1/0/0/-1.5/saddle,
                                     6.7/-0.35/1/-1/-0.35/stable focus,
                                     10.05/0/1/-1/0/centre}{
    \begin{scope}[xshift=\px cm,scale=0.62]
      \draw[gray!60] (-2,0) -- (2,0); \draw[gray!60] (0,-2) -- (0,2);
      \foreach \u in {-1.6,-0.8,...,1.7}{
        \foreach \w in {-1.6,-0.8,...,1.7}{
          \pgfmathsetmacro{\dx}{\ma*\u+\mb*\w}
          \pgfmathsetmacro{\dy}{\mc*\u+\md*\w}
          \pgfmathsetmacro{\nn}{sqrt(\dx*\dx+\dy*\dy)+0.001}
          \draw[->,thin] (\u,\w) -- ++({0.34*\dx/\nn},{0.34*\dy/\nn});
        }}
      \fill (0,0) circle (2.2pt);
      \node[lbl] at (0,-2.45) {\name};
    \end{scope}}
\end{tikzpicture}
\caption{Four of the linear portraits of the table, drawn as vector fields of
$\dot x=Ax$: a stable node ($A=\operatorname{diag}(-1,-2.5)$), a saddle
($\operatorname{diag}(1,-1.5)$), a stable focus (eigenvalues $-0.35\pm j$), and a
centre (eigenvalues $\pm j$).}
\label{fig:portraits}
\end{figure}

\begin{remark}[back to the pendulum]
Linearization tells us which picture a nonlinear system shows \emph{near} each
equilibrium. For the pendulum of Example~\ref{ex:pend}, the hanging equilibrium
has characteristic polynomial $\lambda^{2}+\frac{b}{m\ell^{2}}\lambda+\frac g\ell$:
a stable focus for light friction and a stable node for heavy friction. The upright
equilibrium has $\lambda^{2}+\frac{b}{m\ell^{2}}\lambda-\frac g\ell$, whose roots
are real and of opposite sign: a saddle. One nonlinear portrait contains both
pictures, a spiral sink at $\theta=0$ and a saddle at $\theta=\pi$, stitched
together far from either.
\end{remark}

%% ============================================================== Part III
\partdivider{Part III --- Transfer Functions}

\lecturebreak{Lecture 4}{14 September 2026}

%%%%%%%%%%%%%%%%%%%%%%%%%%%%%%%%%%%%%%%%%%%%%%%%%%%%%%%%%%%%%%%%%%%%%%%%%%%%%%
\section{The Laplace Transform and Transfer Functions}

\subsection{Laplace transform: the part we need}
The (one-sided) Laplace transform of a signal $f(t)$ is
\begin{equation}\label{eq:laplace}
   F(s) = \Lap\{f\}(s) = \int_{0^{-}}^{\infty} f(t)\,e^{-st}\,\d t ,
   \qquad s\in\mathbb{C}.
\end{equation}
The lower limit is written $0^{-}$: the integral starts \emph{before} $t=0$, so
that an impulse sitting at the origin, and the initial state, are counted.

\smallskip
\begin{center}
\renewcommand{\arraystretch}{2.0}
\begin{tabular}{@{}l@{\qquad\qquad}l@{\qquad\qquad}l@{}}
\toprule
$\delta(t) \;\mapsto\; 1$
  & $u_{s}(t) \;\mapsto\; \dfrac{1}{s}$
  & $t\,u_{s}(t) \;\mapsto\; \dfrac{1}{s^{2}}$\\[6pt]
$t^{n}u_{s}(t) \;\mapsto\; \dfrac{n!}{s^{\,n+1}}$
  & $e^{-\alpha t} \;\mapsto\; \dfrac{1}{s+\alpha}$
  & $f(t-a)\,u_{s}(t-a) \;\mapsto\; e^{-as}F(s)$\\[6pt]
$\sin\omega t \;\mapsto\; \dfrac{\omega}{s^{2}+\omega^{2}}$
  & $\cos\omega t \;\mapsto\; \dfrac{s}{s^{2}+\omega^{2}}$
  & $f*g \;\mapsto\; F(s)\,G(s)$\\
\bottomrule
\end{tabular}
\end{center}

\smallskip\noindent
Here $u_{s}$ is the unit step. The pairs are only half of the toolkit; the other
half is the list of \emph{operations} the transform turns into algebra.

\begin{center}\small
\renewcommand{\arraystretch}{1.55}
\begin{tabular}{@{}lll@{}}
\toprule
& time domain & $s$-domain\\
\midrule
linearity        & $a f(t)+b g(t)$            & $aF(s)+bG(s)$\\
differentiation  & $\dot f(t)$                & $sF(s)-f(0^{-})$\\
integration      & $\int_{0}^{t} f(\tau)\,\d\tau$ & $F(s)/s$\\
time shift       & $f(t-a)\,u_s(t-a)$, $a>0$  & $e^{-as}F(s)$\\
frequency shift  & $e^{-\alpha t}f(t)$        & $F(s+\alpha)$\\
multiplication by $t$ & $t\,f(t)$             & $-F'(s)$\\
convolution      & $(f*g)(t)=\int_0^t f(\tau)g(t-\tau)\,\d\tau$ & $F(s)\,G(s)$\\
\bottomrule
\end{tabular}
\end{center}

Of these, the one that does the real work is differentiation:
\begin{equation}\label{eq:diffrule}
   \Lap\{\dot f\} = sF(s) - f(0),
   \qquad
   \Lap\{\ddot f\} = s^{2}F(s) - s f(0) - \dot f(0).
\end{equation}
The one line to remember is its consequence: \emph{if all initial conditions
vanish, differentiation becomes multiplication by $s$} --- and, by the
integration rule, integration becomes division by $s$. That single fact turns a
linear differential equation into an algebraic one, and it is the entire reason
the transfer function exists.

\begin{remark}[reading the table]
Two of the rows are worth seeing as pictures. The frequency shift says that
damping a signal by $e^{-\alpha t}$ slides its transform to the left by $\alpha$
in the $s$-plane: $\sin\omega t$ has poles at $\pm j\omega$, and
$e^{-\alpha t}\sin\omega t$ has poles at $-\alpha\pm j\omega$. The time shift
says that a pure delay of $a$ seconds multiplies by $e^{-as}$ --- a factor that is
not a ratio of polynomials, which is why delays sit awkwardly with everything
built on rational transfer functions.
\end{remark}

\subsection{Initial and final values}\label{sec:fvt}
Two limits let one read the beginning and the end of a signal straight off its
transform, without inverting it.

\begin{proposition}[initial- and final-value theorems]\label{prop:fvt}
Let $f$ have Laplace transform $F$.
\begin{enumerate}[label=(\roman*),itemsep=2pt]
\item If $F$ is strictly proper, then $f(0^{+}) = \lim_{s\to\infty} sF(s)$.
\item If every pole of $sF(s)$ has strictly negative real part, then
$f(t)$ has a limit as $t\to\infty$, and
\[
   \lim_{t\to\infty} f(t) = \lim_{s\to 0} sF(s).
\]
\end{enumerate}
\end{proposition}

Both follow from the differentiation rule: $\int_{0}^{\infty}\dot f(t)e^{-st}\,\d t
= sF(s)-f(0)$, and letting $s\to0$ (the integral becomes $f(\infty)-f(0)$) or
$s\to\infty$ (the integral vanishes) gives the two statements.

\begin{caution}[the hypothesis of the final-value theorem]
The condition on the poles is not decoration. For $F(s)=1/(s-1)$ the formula
gives $\lim_{s\to0}sF(s)=0$, yet $f(t)=e^{t}$ grows without bound. Likewise
$F(s)=\omega/(s^{2}+\omega^{2})$ gives $0$, but $\sin\omega t$ has no limit at
all. The theorem only reports a final value that exists; it cannot decide
whether one does. That question is answered by the poles, and it is the
subject of stability.
\end{caution}

\begin{example}[the steady offset of the held mass]
The feedback law of \S1.3 produced \eqref{eq:massfb},
$m\ddot y+c\dot y+k(y-h)=(\hat m-m)g$. Put $e=y-h$; with $e(0)=\dot e(0)=0$ and a
constant right-hand side $d=(\hat m-m)g$,
\[
   E(s) = \frac{d}{s\,(ms^{2}+cs+k)} ,
   \qquad
   \lim_{t\to\infty} e(t) = \lim_{s\to0} sE(s) = \frac{d}{k} .
\]
The poles of $sE(s)$ are the roots of $ms^{2}+cs+k$, which lie in the left
half-plane for $c,k>0$, so the theorem applies, and it reproduces
\eqref{eq:offset} without solving the differential equation.
\end{example}

\subsection{The transfer function}
\begin{definition}[transfer function]\label{def:tf}
For a linear time-invariant system with input $u(t)$ and output $y(t)$ that is
\emph{initially at rest} (all initial conditions zero), the \emph{transfer
function} is
\[
   P(s) \;=\; \frac{\Lap\{y(t)\}}{\Lap\{u(t)\}} \;=\; \frac{Y(s)}{U(s)} .
\]
\end{definition}

\begin{figure}[htb]
\centering
\begin{tikzpicture}[>=Stealth]
  \node (b) [draw,minimum width=24mm,minimum height=12mm] {$P(s)$};
  \draw[->,thick] ($(b.west)+(-24mm,0)$) node[left] {$u(t)$}
        -- node[above,sig] {$U(s)$} (b.west);
  \draw[->,thick] (b.east)
        -- node[above,sig] {$Y(s)$} ($(b.east)+(24mm,0)$) node[right] {$y(t)$};
\end{tikzpicture}
\caption{The transfer function as a block: $Y(s)=P(s)U(s)$, valid when the system
starts at rest.}
\label{fig:tfblock}
\end{figure}

Equivalently, $P(s)$ is the Laplace transform of the impulse response: putting
$u=\delta$ gives $U\equiv1$ and hence $Y=P$. When several signals are around it
is also written $T_{uy}(s)$, marking which signal is the input and which the
output; for a single-input single-output system the plain $P(s)$ is unambiguous.

\begin{remark}[why ``initially at rest'' is not a detail]
Without that hypothesis the transformed equation carries extra terms
$-sy(0)-\dot y(0)-\cdots$ coming from the initial conditions, and the ratio
$Y(s)/U(s)$ then depends on how the system happened to start, not only on the
system. It would not be a property of the plant, and calling it a transfer
\emph{function of the system} would be meaningless.
\end{remark}

\begin{example}\label{ex:tf}
Take $\ddot y + 3\dot y + y = 2\dot u + u$. Transforming with zero initial
conditions, $s^{2}Y+3sY+Y = 2sU+U$, so
$P(s) = (2s+1)/(s^{2}+3s+1)$.
\end{example}
The recipe is visible in the example: in
$a_n y^{(n)}+\cdots+a_1\dot y+a_0y = b_m u^{(m)}+\cdots+b_1\dot u+b_0u$,
replace every $\d/\d t$ by $s$ and solve for $Y/U$:
\[
   P(s) = \frac{b_ms^{m}+\cdots+b_1s+b_0}{a_ns^{n}+\cdots+a_1s+a_0} .
\]
The degree of the denominator is the order of the system.

\subsection{Example: the two-cart system, idealized}\label{sec:twocarttf}
Strip the two-cart system of friction and of the coupling damper ($c=b=0$):
\begin{equation}\label{eq:twocartideal}
   m_1\ddot p_1 = k\,(p_2-p_1) + u,\qquad
   m_2\ddot p_2 = k\,(p_1-p_2),\qquad y=p_2 .
\end{equation}
Transforming with zero initial conditions, and writing $Y=P_2$,
\[
   m_1 s^{2}P_1 = kY - kP_1 + U,\qquad
   m_2 s^{2}Y = kP_1 - kY .
\]
The second equation gives $P_1 = \dfrac{m_2s^{2}+k}{k}\,Y$. Substituting it into
the first, $\;(m_1s^{2}+k)P_1 = kY+U$, yields
\[
   \frac{(m_1s^{2}+k)(m_2s^{2}+k)}{k}\,Y = kY + U
   \quad\Longrightarrow\quad
   \Bigl[(m_1s^{2}+k)(m_2s^{2}+k) - k^{2}\Bigr]Y = k\,U .
\]
Since $(m_1s^{2}+k)(m_2s^{2}+k)-k^{2} = s^{2}\bigl(m_1m_2s^{2}+k(m_1+m_2)\bigr)$,
\begin{equation}\label{eq:twocartP}
   \boxed{\;
   \frac{Y(s)}{U(s)} \;=\; \frac{k}{\,s^{2}\bigl(m_1m_2\,s^{2}+k(m_1+m_2)\bigr)\,}\;}
\end{equation}

\begin{remark}[reading \eqref{eq:twocartP}]
The factor $s^{2}$ in the denominator is the free rigid-body mode: with no
friction, nothing ties the pair of carts to the ground, so a net impulse leaves
them drifting forever. The remaining factor vanishes at
$s = \pm j\sqrt{k(m_1+m_2)/(m_1m_2)}$, the frequency at which the two carts
oscillate against each other. Both are on the imaginary axis --- with $c=b=0$
nothing damps, which is exactly why the model of Section~2 kept $c$ and $b$.
\end{remark}

\begin{caution}[two different two-cart models]
Equation~\eqref{eq:twocartideal} is \emph{not} the model \eqref{eq:eom}: the
friction $c$ and the damper $b$ have been dropped to keep the algebra readable.
Do not mix results derived from the two. When $c>0$ the factor $s^{2}$ in
\eqref{eq:twocartP} becomes $s(\,\cdot\,)$ --- one pole moves off the origin and
only one free integrator survives.
\end{caution}

\begin{remark}[the same computation with $c$ and $b$ kept]\label{rem:fullTF}
Carrying the full model \eqref{eq:eom} through the same steps gives, in matrix
form,
\[
   \begin{bmatrix}
     m_1s^{2}+(b+c)s+k & -(bs+k)\\[3pt]
     -(bs+k) & m_2s^{2}+(b+c)s+k
   \end{bmatrix}
   \begin{bmatrix}P_1\\[3pt] P_2\end{bmatrix}
   = \begin{bmatrix}U\\[3pt] 0\end{bmatrix},
\]
so that, writing
$\Delta(s) = \bigl[m_1s^{2}+(b+c)s+k\bigr]\bigl[m_2s^{2}+(b+c)s+k\bigr]-(bs+k)^{2}$,
one gets $Y(s)/U(s) = (bs+k)/\Delta(s)$.
Setting $b=c=0$ recovers \eqref{eq:twocartP} exactly, which is a useful check on
both. Two things are worth noticing. The coupling damper produces a finite zero
at $s=-k/b$ --- the numerator is $bs+k$, never $0$. And with $c=0$ but $b>0$ one
gets $\Delta(s)=s^{2}\bigl[m_1m_2s^{2}+(m_1+m_2)bs+(m_1+m_2)k\bigr]$: the double
pole at the origin survives, since it comes from the absence of \emph{ground}
friction, not from the absence of the damper.
\end{remark}

\subsection{From state space to transfer function}\label{sec:ss2tf}
The two representations of Figure~\ref{fig:routes} are connected by a single
computation. Start from \eqref{eq:std} with $x(0)=0$ and transform:
$s\,X(s) = A\,X(s) + B\,U(s)$. Collect the $X$ terms, $\;(sI-A)X(s) = B\,U(s)$,
and invert: $X(s) = (sI-A)^{-1}B\,U(s)$. Feeding this into $Y = CX + DU$,
\begin{equation}\label{eq:ss2tf}
   \boxed{\;Y(s) = \bigl[\,C\,(sI-A)^{-1}B + D\,\bigr]\,U(s),
   \qquad P(s) = C(sI-A)^{-1}B + D. \;}
\end{equation}
For a SISO system this is a scalar rational function of $s$; in general it is a
$p\times m$ matrix of them.

\begin{remark}[keeping $x(0)$, and where the transfer function comes from]
It is worth doing the computation once without assuming $x(0)=0$, because it
shows exactly what the zero-initial-condition hypothesis buys. Transforming
$\dot x = Ax+Bu$ in full gives $sX(s)-x(0) = AX(s)+BU(s)$, hence
$(sI-A)X(s) = x(0) + BU(s)$ and
\[
   Y(s) = \underbrace{C(sI-A)^{-1}x(0)}_{\text{free response}}
        + \underbrace{\bigl[C(sI-A)^{-1}B+D\bigr]U(s)}_{\text{forced response}} .
\]
The transfer function is the second term alone. The first term does not
disappear in reality; it is simply not part of the definition, which is why
Definition~\ref{def:tf} has to say ``initially at rest''.
\end{remark}

\begin{remark}
Since $(sI-A)^{-1} = \operatorname{adj}(sI-A)/\det(sI-A)$, we may also write
$P(s) = C\,\operatorname{adj}(sI-A)\,B/\det(sI-A) + D$, so every pole of $P(s)$
is an eigenvalue of $A$ --- though not necessarily the other way round, since
factors can cancel. The state-space model can therefore carry information the
transfer function does not, which is one reason this course does not work with
transfer functions alone. Going the other way, from $P(s)$ back to a quadruple
$(A,B,C,D)$, is called \emph{realization}; it is always possible for a proper
$P$, but the answer is far from unique. That is why the arrow between the two
boxes in Figure~\ref{fig:routes} points both ways.
\end{remark}

\begin{example}[a realization, step by step]\label{ex:realization}
Realize $P(s)=\dfrac{s^{2}+s+2}{s^{3}+3s^{2}+2s-1}$. Split it through an
intermediate signal $\bar y$:
\[
   U \;\longrightarrow\; \boxed{\dfrac{1}{s^{3}+3s^{2}+2s-1}} \;\longrightarrow\;
   \bar Y \;\longrightarrow\; \boxed{s^{2}+s+2} \;\longrightarrow\; Y .
\]
The first block is the differential equation
$\dddot{\bar y}+3\ddot{\bar y}+2\dot{\bar y}-\bar y=u$, which has no input
derivatives, so the choice $x_1=\bar y$, $x_2=\dot{\bar y}$, $x_3=\ddot{\bar y}$
of \S3.1 turns it into
\[
   \dot x_1=x_2,\qquad \dot x_2=x_3,\qquad \dot x_3 = x_1-2x_2-3x_3+u .
\]
The second block is only differentiation of $\bar y$, and every derivative it needs
is already a state: $y=\ddot{\bar y}+\dot{\bar y}+2\bar y = 2x_1+x_2+x_3$. Hence
\[
   A=\begin{bmatrix}0&1&0\\0&0&1\\1&-2&-3\end{bmatrix},\qquad
   B=\begin{bmatrix}0\\0\\1\end{bmatrix},\qquad
   C=\begin{bmatrix}2&1&1\end{bmatrix},\qquad D=0 .
\]
The pattern is general: the denominator coefficients fill the last row of $A$
(with signs reversed), and the numerator coefficients become $C$. As a check,
$(sI-A)^{-1}B=(1,\;s,\;s^{2})\tp/(s^{3}+3s^{2}+2s-1)$, and $C$ times that is
$P(s)$.
\end{example}

%% ---------------------------------------------------------------- Lecture 5
\lecturebreak{Lecture 5}{16 September 2026}

\subsection{Poles, zeros, and properness}\label{sec:polezero}
Every transfer function we shall meet is a ratio of polynomials,
\begin{equation}\label{eq:rational}
   G(s) = \frac{b_ms^{m}+\cdots+b_1s+b_0}
               {s^{n}+a_{n-1}s^{n-1}+\cdots+a_1s+a_0},
\end{equation}
a \emph{rational} transfer function. Its \emph{degree} is $n$, the degree of the
denominator, and its \emph{relative degree} is $n-m$. Comparing $n$ with $m$ gives
four names:

\begin{center}\small
\begin{tabular}{@{}ll@{}}
\toprule
$n \ge m$ & proper\\
$n > m$   & strictly proper\\
$n = m$   & biproper\\
$n < m$   & improper\\
\bottomrule
\end{tabular}
\end{center}

\begin{remark}[why improper is unwelcome]
An improper transfer function has more derivatives on the output side than the
input side, so realising it requires \emph{differentiating} the input. A pure
differentiator has $|j\omega|=\omega$: its gain grows without bound with
frequency, so it amplifies exactly the high-frequency measurement noise one wants
to suppress. Real plants are strictly proper, which is another way of saying that
nothing responds instantaneously; controllers, which we build ourselves, are kept
at least proper for the same reason.
\end{remark}

Factoring numerator and denominator,
\begin{equation}\label{eq:factored}
   G(s) = \frac{N(s)}{D(s)}
        = k\,\frac{\prod_{i=1}^{m}(s-z_i)}{\prod_{i=1}^{n}(s-p_i)},
\end{equation}
where $N$ and $D$ are real polynomials. The $s\in\mathbb{C}$ with $D(s)=0$ are the
\textbf{poles} --- there are $n$ of them --- and the $s$ with $N(s)=0$ are the
\textbf{zeros}, of which there are $m$. Since $N$ and $D$ have \emph{real}
coefficients, any pole or zero that is not real arrives together with its complex
conjugate. Marking them in the complex plane, poles with $\times$ and zeros with
$\circ$, gives the \emph{pole--zero map} (Figure~\ref{fig:pzmap}).

\begin{figure}[htb]
\centering
\begin{tikzpicture}[>=Stealth,scale=0.95,
   polex/.pic={\draw[thick] (-0.13,-0.13)--(0.13,0.13) (-0.13,0.13)--(0.13,-0.13);}]
  \draw[->] (-3.5,0) -- (1.7,0) node[right] {$\operatorname{Re}s$};
  \draw[->] (0,-2.1) -- (0,2.2) node[above] {$\operatorname{Im}s$};
  \pic at (-2.7,0) {polex};
  \pic at (-1.15,1.35) {polex};
  \pic at (-1.15,-1.35) {polex};
  \draw[thick] (-0.40,0.75) circle (0.14);
  \draw[thick] (-0.40,-0.75) circle (0.14);
  \draw[dashed,gray] (-1.15,1.35) -- (-1.15,-1.35);
  \draw[dashed,gray] (-0.40,0.75) -- (-0.40,-0.75);
  \node[anchor=west,lbl] at (0.30,1.80) {$\times$ pole \qquad $\circ$ zero};
\end{tikzpicture}
\caption{A pole--zero map. Because $N$ and $D$ have real coefficients, every
non-real pole and zero is paired with its mirror image across the real axis
(dashed); real ones stand alone.}
\label{fig:pzmap}
\end{figure}

\subsection{Inverting by partial fractions}\label{sec:pfe}
To go back from $Y(s)$ to $y(t)$ one almost never evaluates an inverse integral.
Instead one splits a rational $Y(s)$ into a sum of simple terms, each of which is
a line of the table of pairs. If $Y$ is strictly proper with poles $p_i$, there
are three kinds of term.

\smallskip\noindent\textbf{Distinct real poles.} Each simple pole $p_i$
contributes $A_i/(s-p_i)$, and its coefficient is found by the \emph{cover-up
rule}: multiply by $(s-p_i)$ and set $s=p_i$,
\begin{equation}\label{eq:coverup}
   A_i = \bigl[(s-p_i)\,Y(s)\bigr]_{s=p_i} .
\end{equation}
Each such term inverts to $A_i e^{p_i t}$.

\smallskip\noindent\textbf{A repeated pole.} A pole of multiplicity $r$
contributes $r$ terms, $B_1/(s-p)+B_2/(s-p)^{2}+\dots+B_r/(s-p)^{r}$. The top
coefficient $B_r$ is found by cover-up with $(s-p)^{r}$; the lower ones by
differentiating: $B_{r-k} = \frac{1}{k!}\frac{\d^{k}}{\d s^{k}}
\bigl[(s-p)^{r}Y(s)\bigr]_{s=p}$. The term $B/(s-p)^{j}$ inverts to
$B\,t^{j-1}e^{pt}/(j-1)!$ --- a repeated pole brings a polynomial in $t$ in front
of the exponential.

\smallskip\noindent\textbf{A complex pair.} Poles $-\alpha\pm j\omega$ are kept
together as one real quadratic, $(s+\alpha)^{2}+\omega^{2}$, over a linear
numerator. Completing it into the two shifted pairs of the table,
\[
   \frac{C(s+\alpha)+D\omega}{(s+\alpha)^{2}+\omega^{2}}
   \;\longmapsto\;
   e^{-\alpha t}\bigl(C\cos\omega t + D\sin\omega t\bigr),
\]
gives a decaying (for $\alpha>0$) oscillation. The real part of the poles sets the
envelope, the imaginary part the frequency.

\begin{example}[the three cases]\label{ex:pfe}
\begin{enumerate}[label=(\alph*),itemsep=3pt,leftmargin=*]
\item $Y(s)=\dfrac{s+5}{(s+1)(s+2)}=\dfrac{A}{s+1}+\dfrac{B}{s+2}$. Cover-up gives
$A=\frac{s+5}{s+2}\big|_{s=-1}=4$ and $B=\frac{s+5}{s+1}\big|_{s=-2}=-3$, so
$y(t)=4e^{-t}-3e^{-2t}$.
\item $Y(s)=\dfrac{1}{s(s+1)^{2}}=\dfrac{A}{s}+\dfrac{B_1}{s+1}+\dfrac{B_2}{(s+1)^{2}}$.
Cover-up gives $A=1$ and $B_2=\frac1s\big|_{s=-1}=-1$; differentiating,
$B_1=\frac{\d}{\d s}\frac1s\big|_{s=-1}=-1$. Hence
$y(t)=1-e^{-t}-t\,e^{-t}$.
\item $Y(s)=\dfrac{1}{s^{2}+2s+5}=\dfrac{1}{(s+1)^{2}+2^{2}}
=\dfrac12\cdot\dfrac{2}{(s+1)^{2}+2^{2}}$, so $y(t)=\tfrac12e^{-t}\sin 2t$.
\end{enumerate}
In (b) the two theorems of \S\ref{sec:fvt} check the answer:
$\lim_{s\to\infty}sY(s)=0=y(0^{+})$ and $\lim_{s\to0}sY(s)=1=y(\infty)$.
\end{example}

\subsection{Where the poles of the output come from}\label{sec:outputpoles}
With the plant in hand, the response to an input is
$y(t) = \Lap^{-1}\{G(s)\,U(s)\}$, and the way to read that inverse transform is to
split it into partial fractions. Suppose it comes out as
\begin{equation}\label{eq:pfe}
   Y(s) = \frac{A}{s-1} + \frac{B}{s-2} + \frac{C}{s+3}
        + \frac{Ds+E}{s^{2}+\omega^{2}} .
\end{equation}
The first three denominators are poles of the \emph{plant}; the last is a pole of
the \emph{input} --- an input $\sin\omega t$ would contribute exactly that factor.
Inverting term by term gives one time function per pole, so
\begin{center}
\emph{the poles of the output are the poles of the plant together with the poles
of the input.}
\end{center}
The plant's own poles appear whatever we drive it with, which is why the question
of whether they are well placed is a question about the plant alone.

\subsection{Block diagrams and the feedback loop}\label{sec:blocks}
Because $Y(s)=P(s)U(s)$ is plain multiplication, systems wired together combine by
algebra. Three connections cover almost every diagram (Figure~\ref{fig:blocks}):
\begin{equation}\label{eq:blockalg}
   \underbrace{\;Q(s)P(s)\;}_{\text{series}},
   \qquad
   \underbrace{\;P(s)+Q(s)\;}_{\text{parallel}},
   \qquad
   \underbrace{\;\frac{G(s)}{1+G(s)H(s)}\;}_{\text{negative feedback}} .
\end{equation}
The first two are immediate. The third is the one that matters, and it is worth
deriving once. In the loop of Figure~\ref{fig:blocks}(c) the error is
$E=R-HY$ and the output is $Y=GE$, so $Y=G(R-HY)$: the output appears on both
sides, which is the algebraic trace of feedback. Collecting it,
$(1+GH)Y=GR$, and the formula follows.

The trade-off against the equations of motion of \S2 is worth stating. A block
diagram makes \emph{connecting} systems trivial --- one multiplies, adds, or applies
\eqref{eq:blockalg} --- but it hides the physics: a block labelled $P(s)$ no longer
says which masses, springs and dampers are inside it. The differential equation is
the opposite: every term has a physical meaning, and wiring two systems together
means rewriting the equations. The two are used side by side for that reason.

\begin{figure}[htb]
\centering
\begin{tikzpicture}[>=Stealth,node distance=9mm]
  % (a) series
  \node[lbl] (la) at (-1.1,0) {(a)};
  \node (ua) at (0,0) {$U$};
  \node[draw,minimum width=10mm,minimum height=7mm] (pa) at (1.5,0) {$P$};
  \node[draw,minimum width=10mm,minimum height=7mm] (qa) at (3.2,0) {$Q$};
  \node (ya) at (4.7,0) {$Y$};
  \draw[->] (ua) -- (pa); \draw[->] (pa) -- (qa); \draw[->] (qa) -- (ya);
  % (b) parallel
  \node[lbl] at (6.0,0) {(b)};
  \node (ub) at (6.8,0) {$U$};
  \coordinate (sb) at (7.4,0);
  \node[draw,minimum width=10mm,minimum height=7mm] (pb) at (8.6,0.55) {$P$};
  \node[draw,minimum width=10mm,minimum height=7mm] (qb) at (8.6,-0.55) {$Q$};
  \node[sumc] (jb) at (9.9,0) {};
  \node (yb) at (10.9,0) {$Y$};
  \draw (ub) -- (sb); \fill (sb) circle (1.2pt);
  \draw[->] (sb) |- (pb); \draw[->] (sb) |- (qb);
  \draw[->] (pb) -| node[pos=0.8,right,lbl] {$+$} (jb);
  \draw[->] (qb) -| node[pos=0.8,right,lbl] {$+$} (jb);
  \draw[->] (jb) -- (yb);
  % (c) feedback
  \node[lbl] at (1.2,-2.2) {(c)};
  \node (rc) at (2.0,-2.2) {$R$};
  \node[sumc] (sc) at (3.1,-2.2) {};
  \node[draw,minimum width=10mm,minimum height=7mm] (gc) at (4.9,-2.2) {$G$};
  \coordinate (tc) at (6.5,-2.2);
  \node (yc) at (7.4,-2.2) {$Y$};
  \node[draw,minimum width=10mm,minimum height=7mm] (hc) at (4.9,-3.3) {$H$};
  \draw[->] (rc) -- node[pos=0.75,above,lbl] {$+$} (sc);
  \draw[->] (sc) -- node[above,sig] {$E$} (gc);
  \draw[->] (gc) -- (yc);
  \fill (tc) circle (1.2pt);
  \draw[->] (tc) |- (hc);
  \draw[->] (hc) -| node[pos=0.85,left,lbl] {$-$} (sc);
\end{tikzpicture}
\caption{(a) Series: $Y=QPU$. (b) Parallel: $Y=(P+Q)U$. (c) Negative feedback:
$Y=\frac{G}{1+GH}R$.}
\label{fig:blocks}
\end{figure}

In the standard control loop the forward path is a controller followed by the
plant, $G=P(s)C(s)$, and the output is fed back directly, $H=1$. The reference
$r$ is what we want $y$ to be, and the controller acts on the error $e=r-y$:
\begin{equation}\label{eq:closedloop}
   \frac{Y(s)}{R(s)} = \frac{P(s)C(s)}{1+P(s)C(s)},
   \qquad
   \frac{E(s)}{R(s)} = \frac{1}{1+P(s)C(s)} .
\end{equation}

\begin{definition}[loop transfer function]\label{def:loop}
The product of all the blocks met in going once around the loop,
\[
   L(s) = P(s)\,C(s)
\]
(more generally $G(s)H(s)$), is the \emph{loop transfer function} --- also called
the \emph{loop gain}, or the \emph{open-loop} transfer function. The equation
\begin{equation}\label{eq:chareq}
   1+L(s)=0
\end{equation}
is the \emph{characteristic equation} of the closed loop.
\end{definition}

Every closed-loop transfer function in \eqref{eq:closedloop} has the same
denominator $1+L$. Writing $L=N_L/D_L$ as a ratio of polynomials,
\[
   \frac{L}{1+L} = \frac{N_L}{D_L+N_L},
\]
so the closed-loop poles are the roots of $D_L+N_L$, equivalently of
\eqref{eq:chareq}. This is the precise sense in which \emph{feedback moves poles}:
the plant alone has the roots of $D_L$; closing the loop replaces them by the
roots of $D_L+N_L$, and the controller $C$ is our handle on the difference.

\begin{example}[the held mass, as a loop]\label{ex:massloop}
Go back to \S1.3 and measure everything from the target height, $e=h-y$. The
plant from force to height is $P(s)=1/(ms^{2})$, and the feedback law
\eqref{eq:pd} is, apart from the gravity-compensating constant, the controller
$C(s)=k+cs$. Then
\[
   L(s) = \frac{k+cs}{ms^{2}},
   \qquad
   1+L(s) = \frac{ms^{2}+cs+k}{ms^{2}} .
\]
The open loop has a double pole at the origin --- a free mass, which drifts. The
closed loop has the roots of $ms^{2}+cs+k$, which for $c,k>0$ lie in the left
half-plane. The gains have become an \emph{effective damping} $c$ and an
\emph{effective stiffness} $k$ that the free mass did not have, and that is the
whole content of ``feedback stabilised it''.
\end{example}

\begin{question}
Given a plant $P$ and a proposed controller $C$, is the closed loop stable? One
answer is to compute the roots of $D_L+N_L$. Section~\ref{sec:nyquist} gives a
second answer, which needs only the graph of $L(j\omega)$ --- something one can
measure --- and which tells us not just \emph{whether} the loop is stable but
\emph{how far} it is from losing stability.
\end{question}

%% =============================================================== Part IV
\partdivider{Part IV --- Frequency Response and Stability}

%%%%%%%%%%%%%%%%%%%%%%%%%%%%%%%%%%%%%%%%%%%%%%%%%%%%%%%%%%%%%%%%%%%%%%%%%%%%%%
\section{Frequency Response and Bode Plots}\label{sec:freq}

\subsection{Stable plants}
\begin{definition}[stable plant]\label{def:stable}
A plant is \emph{stable} when every pole has strictly negative real part.
\end{definition}

\begin{remark}[why this section begins with stability]
Everything below evaluates $P$ on the imaginary axis, at $s=j\omega$ --- that is,
it trades the Laplace transform for the Fourier transform --- and that substitution
is only legitimate when the Laplace integral actually converges there. For a
stable plant it does, and $P(j\omega)$ is the Fourier transform of the impulse
response --- something one can measure by driving the plant with a sinusoid. If a
pole sits in the right half-plane the imaginary axis lies outside the region of
convergence, the response to a sinusoid grows without bound, and the symbol
$P(j\omega)$ is not the frequency response of anything. Stability is a hypothesis
of this whole section, not a side remark.
\end{remark}

\subsection{Sinusoid in, sinusoid out}
Drive a stable plant with $\sin\omega t$ and wait for the transient to die away.
What is left is a sinusoid \emph{at the same frequency}:
\begin{equation}\label{eq:freqresp}
   \sin\omega t \;\longmapsto\; A\sin(\omega t+\theta),
   \qquad A = |P(j\omega)|,\qquad \theta = \angle P(j\omega).
\end{equation}
Only two numbers change, and both are read off the single complex number
$P(j\omega)$: its magnitude scales the amplitude, its argument shifts the phase.
As functions of $\omega$ the pair is the \emph{frequency response}. This is the
practical reason transfer functions are worth having: it turns a differential
equation into a gain and a phase shift, one frequency at a time.

\begin{proof}[Why \eqref{eq:freqresp} holds]
The input transform is $\omega/(s^{2}+\omega^{2})$, so by \S\ref{sec:outputpoles}
\[
   Y(s) = P(s)\,\frac{\omega}{s^{2}+\omega^{2}}
        = \frac{as+b}{s^{2}+\omega^{2}} + \bigl(\text{terms from the poles of }P\bigr).
\]
The plant is stable, so the second group dies out, and what remains is
$y_{ss}(t) = a\cos\omega t+(b/\omega)\sin\omega t$. To find $a$ and $b$, multiply
both sides by $s-j\omega$ and set $s=j\omega$; the plant terms vanish and
\[
   P(j\omega)\,\frac{\omega}{2j\omega} = \frac{a\,j\omega+b}{2j\omega}
   \quad\Longrightarrow\quad
   P(j\omega) = \frac{b}{\omega} + j a .
\]
So $b/\omega=|P(j\omega)|\cos\angle P(j\omega)$ and $a=|P(j\omega)|\sin\angle
P(j\omega)$, and the sum-of-angles formula collapses $y_{ss}$ to
$|P(j\omega)|\sin\bigl(\omega t+\angle P(j\omega)\bigr)$.
\end{proof}

\begin{example}[two frequencies at once]
Feed $u(t)=2\sin t+3\cos(2t+1)$ into $P(s)=1/(s+1)$. By linearity each term is
handled on its own, and a cosine obeys the same rule as a sine. At $\omega=1$,
$P(j)=1/(1+j)$ has magnitude $1/\sqrt2$ and phase $-\pi/4$; at $\omega=2$,
$P(2j)=1/(1+2j)$ has magnitude $1/\sqrt5$ and phase $-\arctan2\approx-63.4^\circ$.
Hence
\[
   y_{ss}(t) = \sqrt2\,\sin\!\Bigl(t-\frac\pi4\Bigr)
             + \frac{3}{\sqrt5}\cos\bigl(2t+1-\arctan 2\bigr).
\]
The higher frequency is both attenuated more and delayed more: $1/(s+1)$ is a
low-pass filter, which the Bode plot of the next subsection shows at a glance.
\end{example}

\subsection{The Bode plot}
The frequency response is drawn as two graphs against $\log_{10}\omega$: the
magnitude in decibels,
\[
   20\log_{10}|P(j\omega)| ,
\]
and the phase $\angle P(j\omega)$. Together they are the \emph{Bode plot}.

Both logarithms earn their place. A transfer function is a \emph{product} of
factors, and taking the logarithm of the magnitude turns that product into a
\emph{sum}, so each factor contributes its own piece and the whole is obtained by
adding the pieces up. Putting $\log\omega$ on the horizontal axis then makes each
of those pieces a straight line over most of its range.

\subsection{Building the magnitude from the factors}
Take, for example,
\begin{equation}\label{eq:bodeex}
   P(s) = \frac{K(s+z_1)(s+z_2)}
               {(s+p_1)(s+p_2)^{2}\bigl(s^{2}+2\zeta\omega_n s+\omega_n^{2}\bigr)} .
\end{equation}
Taking $20\log_{10}$ of the magnitude turns every factor into a term:
\begin{equation}\label{eq:bodesum}
\begin{aligned}
   20\log|P| = 20\bigl(&\log|s+z_1| + \log|s+z_2| - \log|s+p_1| \\
                       &- 2\log|s+p_2| - \log|s^{2}+2\zeta\omega_n s+\omega_n^{2}|\bigr),
\end{aligned}
\end{equation}
evaluated at $s=j\omega$. Numerator factors enter with a plus, denominator factors
with a minus, and a repeated factor enters with its multiplicity --- hence the
$2$ in front of $\log|s+p_2|$.

Each first-order factor is then a single bend (Figure~\ref{fig:bode}). For
$s+z_1$ at $s=j\omega$,
\[
   20\log|j\omega+z_1| = 10\log_{10}\bigl(\omega^{2}+z_1^{2}\bigr)
   \;\approx\;
   \begin{cases}
      20\log z_1, & \omega \ll z_1,\\[2pt]
      20\log \omega, & \omega \gg z_1,
   \end{cases}
\]
so the asymptotes are a horizontal line and a line of slope $+20$\,dB per decade,
meeting at the \emph{corner frequency} $\omega=z_1$. At the corner itself the true
value is $20\log(z_1\sqrt2) = 20\log z_1 + 3.01$\,dB: the curve passes $3$\,dB
\emph{above} the corner. A denominator factor gives the same picture upside down
--- slope $-20$\,dB/decade and $3$\,dB \emph{below} --- and the repeated factor
$(s+p_2)^{2}$ doubles both, to $-40$\,dB/decade.

\begin{figure}[htb]
\centering
\begin{tikzpicture}[>=Stealth,x=1.75cm,y=0.085cm]
  % ---- axes ----
  \draw[->] (-1.62,0) -- (1.80,0) node[below right,lbl] {$\log_{10}\omega$};
  \draw[->] (-1.5,-6) -- (-1.5,35);
  \node[rotate=90,anchor=south,lbl] at (-1.78,16) {$20\log_{10}|j\omega+z_1|$};
  % ---- corner ----
  \draw[dashed,gray] (0,-6) -- (0,33);
  \node[below,lbl] at (0,-1.5) {$\log_{10}z_1$};
  % ---- true curve ----
  \draw[gray!65,thick] plot[domain=-1.5:1.5,samples=90,smooth,variable=\u]
        (\u,{10*log10(1+pow(10,2*\u))});
  % ---- asymptotes ----
  \draw[thick] (-1.5,0) -- (0,0);
  \draw[thick] (0,0) -- (1.5,30);
  % ---- 3 dB at the corner ----
  \node[circle,fill,inner sep=1.1pt] at (0,3.01) {};
  \node[left,lbl] at (-0.05,7) {$3$\,dB};
  \node[above left,lbl] at (1.28,29) {$+20$\,dB/decade};
\end{tikzpicture}
\caption{One first-order factor. The two straight asymptotes meet at the corner
frequency $\omega=z_1$; the true magnitude (grey) runs $3$\,dB above the corner.
A Bode magnitude plot is built by adding one such picture per factor.}
\label{fig:bode}
\end{figure}

\subsection{Phase by adding angles}\label{sec:phase}
The phase is even easier, because the same factorisation is already a product of
complex numbers. Writing each factor in polar form at $s=j\omega$,
\[
   P(j\omega)
   = \frac{\prod_i M_i e^{\,j\theta_i}}{\prod_j M'_j e^{\,j\theta'_j}} ,
   \qquad\text{so}\qquad
   \angle P(j\omega) = \sum_i \theta_i - \sum_j \theta'_j .
\]
Each factor's angle is measured separately and they are simply added and
subtracted --- the phase plot, like the magnitude plot, is assembled one factor at
a time. For a first-order factor the angle is
\[
   \angle(j\omega+z_1) = \arctan\frac{\omega}{z_1},
\]
which climbs from $0^\circ$ at low frequency to $90^\circ$ at high frequency and
passes $45^\circ$ exactly at the corner. Almost all of the change happens within a
decade on either side, so the usual straight-line sketch holds $0^\circ$ up to
$z_1/10$, rises linearly in $\log\omega$ to $90^\circ$ at $10z_1$, and stays there.
A denominator factor contributes the negative of this.

\begin{remark}[a geometric reading]\label{rem:vectors}
$j\omega+z_1 = j\omega-(-z_1)$ is the vector from the zero at $-z_1$ to the point
$j\omega$ on the imaginary axis. As $\omega$ runs up the axis, its length is the
factor's magnitude and its angle is the factor's phase. Every factor of
$P(j\omega)$ is such a vector, drawn from a pole or a zero to $j\omega$: the
magnitude of $P$ is the product of the zero-vector lengths divided by the product
of the pole-vector lengths, and the phase is the sum of the zero-vector angles
minus the sum of the pole-vector angles. The same picture, with $j\omega$ replaced
by a point running around a closed curve, is the key to
Section~\ref{sec:nyquist}.
\end{remark}

%% ---------------------------------------------------------------- Lecture 6
\lecturebreak{Lecture 6}{21 September 2026}

\subsection{The second-order system and the damping ratio}\label{sec:second}
The quadratic factor in \eqref{eq:bodeex} deserves its own treatment: it is the one
factor whose Bode plot the asymptotes do \emph{not} describe well, and it is the
transfer function of every damped oscillator. The standard second-order system is
\begin{equation}\label{eq:second}
   G(s) = \frac{\omega_n^{2}}{s^{2}+2\zeta\omega_n s+\omega_n^{2}},
   \qquad \omega_n>0,\quad \zeta\ge 0,
\end{equation}
normalised so that $G(0)=1$. The two parameters have names: $\omega_n$ is the
\emph{natural frequency} and $\zeta$ the \emph{damping ratio}. Both are forced on
us by physics. A mass on a spring with a damper, $m\ddot q+b\dot q+kq=ku$, divided
through by $m$ is exactly \eqref{eq:second} with
\begin{equation}\label{eq:zetawn}
   \omega_n = \sqrt{k/m},\qquad
   \zeta = \frac{b}{2\sqrt{km}} ,
\end{equation}
so $\omega_n$ is the frequency at which the system would oscillate if there were no
damping, and $\zeta$ measures the actual damping against the amount $2\sqrt{km}$
that just suppresses oscillation. Because the damping term is written
$2\zeta\omega_n s$ and not $2\zeta s$, $\zeta$ is a pure number, and everything
said below about $\zeta$ holds whatever the time scale of the system.

\smallskip\noindent\textbf{Poles.} The roots of $s^{2}+2\zeta\omega_n s+\omega_n^{2}$ are
\[
   s = -\zeta\omega_n \pm \omega_n\sqrt{\zeta^{2}-1},
\]
and the value of $\zeta$ sorts the system into four regimes:
\begin{center}\small
\renewcommand{\arraystretch}{1.3}
\begin{tabular}{@{}lll@{}}
\toprule
$\zeta>1$   & overdamped        & two distinct negative real poles\\
$\zeta=1$   & critically damped & a double pole at $-\omega_n$\\
$0<\zeta<1$ & underdamped       & $-\zeta\omega_n\pm j\omega_d$, \quad $\omega_d=\omega_n\sqrt{1-\zeta^{2}}$\\
$\zeta=0$   & undamped          & $\pm j\omega_n$, on the imaginary axis\\
\bottomrule
\end{tabular}
\end{center}
In the underdamped case, which is the interesting one, the poles lie on the circle
$|s|=\omega_n$, and the line from the origin to either pole makes an angle
$\theta$ with the negative real axis given by
\begin{equation}\label{eq:zetaangle}
   \cos\theta = \zeta
\end{equation}
(Figure~\ref{fig:secondpoles}). Distance from the origin is frequency; the angle is
damping. Lines of constant $\zeta$ are rays from the origin.

\begin{figure}[htb]
\centering
\begin{tikzpicture}[>=Stealth,scale=1.0,
   polex/.pic={\draw[thick] (-0.12,-0.12)--(0.12,0.12) (-0.12,0.12)--(0.12,-0.12);}]
  \draw[->] (-2.9,0) -- (1.3,0) node[right] {$\operatorname{Re}s$};
  \draw[->] (0,-2.4) -- (0,2.5) node[above] {$\operatorname{Im}s$};
  \draw[dashed,gray] (90:2) arc (90:270:2);
  \draw[thin] (0,0) -- (-1,1.732);
  \draw[thin] (0,0) -- (-1,-1.732);
  \pic at (-1,1.732) {polex};
  \pic at (-1,-1.732) {polex};
  \draw[dotted] (-1,1.732) -- (-1,0);
  \draw[dotted] (-1,1.732) -- (0,1.732);
  \node[below,lbl] at (-1,-0.02) {$-\zeta\omega_n$};
  \node[right,lbl] at (0.02,1.732) {$j\omega_d$};
  \draw (-0.6,0) arc (180:120:0.6);
  \node[lbl] at (-0.78,0.33) {$\theta$};
  \node[lbl,above left] at (-0.48,0.86) {$\omega_n$};
  \node[lbl,anchor=west] at (0.35,-1.25) {$\cos\theta=\zeta$};
\end{tikzpicture}
\caption{The underdamped pole pair of \eqref{eq:second}, drawn for $\zeta=0.5$.
The poles sit on the circle of radius $\omega_n$; the angle $\theta$ they make with
the negative real axis is fixed by the damping ratio alone.}
\label{fig:secondpoles}
\end{figure}

\smallskip\noindent\textbf{Time response.} For $0<\zeta<1$, the partial fractions
of \S\ref{sec:pfe} applied to the unit step, $Y(s)=G(s)/s$, give
\begin{equation}\label{eq:secondstep}
   y(t) = 1 - \frac{e^{-\zeta\omega_n t}}{\sqrt{1-\zeta^{2}}}\,
          \sin\bigl(\omega_d t+\theta\bigr),
   \qquad \theta=\arccos\zeta .
\end{equation}
Every piece of this is read off the pole: the real part $-\zeta\omega_n$ sets the
decaying envelope, the imaginary part $\omega_d$ the frequency of the ringing
(Figure~\ref{fig:secondstep}). Small $\zeta$ gives a quick but badly ringing
response; $\zeta$ near $1$ a sluggish one that barely overshoots. Two numbers
summarise the curve, the peak overshoot and the settling time:
\begin{equation}\label{eq:overshoot}
   M_p = e^{-\pi\zeta/\sqrt{1-\zeta^{2}}},
   \qquad
   t_s \approx \frac{4}{\zeta\omega_n}\quad\text{(to within $2\%$)} .
\end{equation}
The overshoot depends on $\zeta$ alone, that is on the \emph{angle} of the poles;
the settling time on $\zeta\omega_n$ alone, that is on their \emph{distance from
the imaginary axis}.

\begin{figure}[htb]
\centering
\begin{tikzpicture}[>=Stealth,x=0.46cm,y=2.1cm]
  \draw[->] (0,0) -- (14.7,0) node[right,lbl] {$\omega_n t$};
  \draw[->] (0,0) -- (0,1.95) node[above,lbl] {$y(t)$};
  \draw[dotted,gray] (0,1) -- (14,1);
  \foreach \tk in {0,2,...,14}{\draw (\tk,0) -- (\tk,-0.04) node[below,lbl] {\tk};}
  \foreach \yk in {0.5,1,1.5}{\draw (0,\yk) -- (-0.18,\yk) node[left,lbl] {\yk};}
  \draw[gray] plot[domain=0:14,samples=150,smooth]
    (\x,{1-exp(-0.1*\x)/sqrt(0.99)*sin(deg(sqrt(0.99)*\x)+acos(0.1))});
  \draw[dashed] plot[domain=0:14,samples=120,smooth]
    (\x,{1-exp(-0.3*\x)/sqrt(0.91)*sin(deg(sqrt(0.91)*\x)+acos(0.3))});
  \draw[thick] plot[domain=0:14,samples=100,smooth]
    (\x,{1-exp(-0.7*\x)/sqrt(0.51)*sin(deg(sqrt(0.51)*\x)+acos(0.7))});
  \draw[thick,dotted] plot[domain=0:14,samples=80,smooth] (\x,{1-(1+\x)*exp(-\x)});
  % legend
  \begin{scope}[shift={(9.2,0.14)}]
    \draw[gray] (0,0.36) -- (1.4,0.36); \node[right,lbl] at (1.5,0.36) {$\zeta=0.1$};
    \draw[dashed] (0,0.24) -- (1.4,0.24); \node[right,lbl] at (1.5,0.24) {$\zeta=0.3$};
    \draw[thick] (0,0.12) -- (1.4,0.12); \node[right,lbl] at (1.5,0.12) {$\zeta=0.7$};
    \draw[thick,dotted] (0,0) -- (1.4,0); \node[right,lbl] at (1.5,0) {$\zeta=1$};
  \end{scope}
\end{tikzpicture}
\caption{Unit-step responses of \eqref{eq:second} against $\omega_n t$. Less damping
means a higher first peak and slower-dying ringing; $\zeta=1$ does not overshoot at
all.}
\label{fig:secondstep}
\end{figure}

\begin{remark}[the semester's goal, restated]
Question~\ref{q:goal} asks the carts to settle within $5$ seconds. By
\eqref{eq:overshoot}, how fast a response settles and how far it overshoots are
both statements about \emph{where the dominant poles are}: settling within $5$\,s
asks for $\zeta\omega_n \gtrsim 4/5$, and the overshoot is fixed by the angle
$\theta$. The open-loop two-cart system cannot be told where to put its poles.
A closed loop can --- that is what \S\ref{sec:blocks} showed feedback doing.
\end{remark}

\smallskip\noindent\textbf{Frequency response.} On the imaginary axis,
\begin{equation}\label{eq:secondfreq}
   G(j\omega) = \frac{\omega_n^{2}}{(\omega_n^{2}-\omega^{2})+j\,2\zeta\omega_n\omega} .
\end{equation}
With the normalised frequency $u=\omega/\omega_n$,
\[
   |G(j\omega)| = \frac{1}{\sqrt{(1-u^{2})^{2}+(2\zeta u)^{2}}},
   \qquad
   \angle G(j\omega) = -\operatorname{atan2}\bigl(2\zeta u,\;1-u^{2}\bigr).
\]
The asymptotes are what one expects of a quadratic: for $u\ll1$ the magnitude is
$1$ (0\,dB) and the phase $0^\circ$; for $u\gg1$ the magnitude is $1/u^{2}$, a line
of slope $-40$\,dB/decade, and the phase is $-180^\circ$. They meet at
$\omega=\omega_n$. What happens \emph{at} the corner, however, is decided by
$\zeta$. Setting $u=1$ in \eqref{eq:secondfreq},
\[
   G(j\omega_n) = \frac{1}{j\,2\zeta},
   \qquad
   |G(j\omega_n)| = \frac{1}{2\zeta},
   \qquad
   \angle G(j\omega_n) = -90^\circ .
\]
The phase at the corner is $-90^\circ$ for every $\zeta$, but the magnitude is
$1/(2\zeta)$, which is large when the damping is small: $\zeta=0.1$ puts the curve
$14$\,dB above the asymptotes (Figure~\ref{fig:secondbode}). A first-order factor
never strays more than $3$\,dB from its asymptotes; the quadratic can carry a tall
\emph{resonant peak}. Setting the derivative of $|G|^{2}$ with respect to $u$ to
zero shows that for $\zeta<1/\sqrt2$ the peak lies at
\begin{equation}\label{eq:resonance}
   \omega_r = \omega_n\sqrt{1-2\zeta^{2}},
   \qquad
   M_r = |G(j\omega_r)| = \frac{1}{2\zeta\sqrt{1-\zeta^{2}}},
\end{equation}
and that for $\zeta\ge1/\sqrt2\approx0.707$ there is no peak at all. The phase,
meanwhile, falls from $0^\circ$ to $-180^\circ$ more and more abruptly as $\zeta$
shrinks; in the limit $\zeta\to0$ it jumps by $180^\circ$ at $\omega_n$.

\begin{figure}[htb]
\centering
\begin{tikzpicture}[>=Stealth,x=2.7cm]
  % ---------------- magnitude ----------------
  \begin{scope}[y=0.058cm]
    \draw[->] (-1,0) -- (1.13,0) node[right,lbl] {$\log_{10}(\omega/\omega_n)$};
    \draw[->] (-1,-42) -- (-1,20) node[above,lbl] {dB};
    \foreach \d in {-40,-20,0}{\draw (-1,\d) -- (-1.03,\d) node[left,lbl] {\d};}
    \draw[dotted,gray] (0,-42) -- (0,18);
    \draw[gray!70,densely dashed] (-1,0) -- (0,0) -- (1,-40);
    \draw[gray] plot[domain=-1:1,samples=160,smooth]
      (\x,{-10*log10(pow(1-pow(10,2*\x),2)+pow(0.2*pow(10,\x),2))});
    \draw[dashed] plot[domain=-1:1,samples=120,smooth]
      (\x,{-10*log10(pow(1-pow(10,2*\x),2)+pow(0.6*pow(10,\x),2))});
    \draw[thick] plot[domain=-1:1,samples=100,smooth]
      (\x,{-10*log10(pow(1-pow(10,2*\x),2)+pow(1.4*pow(10,\x),2))});
    \draw[thick,dotted] plot[domain=-1:1,samples=100,smooth]
      (\x,{-10*log10(pow(1-pow(10,2*\x),2)+pow(2.0*pow(10,\x),2))});
    \begin{scope}[shift={(-0.93,-33)}]
      \draw[gray] (0,9) -- (0.2,9); \node[right,lbl] at (0.22,9) {$\zeta=0.1$};
      \draw[dashed] (0,4.5) -- (0.2,4.5); \node[right,lbl] at (0.22,4.5) {$\zeta=0.3$};
      \draw[thick] (0.78,9) -- (0.98,9); \node[right,lbl] at (1.0,9) {$\zeta=0.7$};
      \draw[thick,dotted] (0.78,4.5) -- (0.98,4.5); \node[right,lbl] at (1.0,4.5) {$\zeta=1$};
    \end{scope}
  \end{scope}
  % ---------------- phase ----------------
  \begin{scope}[yshift=-4.3cm,y=0.017cm]
    \draw[->] (-1,0) -- (1.13,0) node[right,lbl] {$\log_{10}(\omega/\omega_n)$};
    \draw[->] (-1,-190) -- (-1,12) node[above,lbl] {phase};
    \foreach \d in {0,-90,-180}{\draw (-1,\d) -- (-1.03,\d) node[left,lbl] {$\d^\circ$};}
    \draw[dotted,gray] (0,-190) -- (0,5);
    \draw[dotted,gray] (-1,-90) -- (1,-90);
    \draw[gray] plot[domain=-1:1,samples=160,smooth]
      (\x,{-atan2(0.2*pow(10,\x),1-pow(10,2*\x))});
    \draw[dashed] plot[domain=-1:1,samples=120,smooth]
      (\x,{-atan2(0.6*pow(10,\x),1-pow(10,2*\x))});
    \draw[thick] plot[domain=-1:1,samples=100,smooth]
      (\x,{-atan2(1.4*pow(10,\x),1-pow(10,2*\x))});
    \draw[thick,dotted] plot[domain=-1:1,samples=100,smooth]
      (\x,{-atan2(2.0*pow(10,\x),1-pow(10,2*\x))});
  \end{scope}
\end{tikzpicture}
\caption{Bode magnitude (top) and phase (bottom) of the second-order system
\eqref{eq:second}. The grey dashed lines in the top panel are the asymptotes. Every
curve passes $-90^\circ$ at $\omega=\omega_n$; the peak there,
$20\log_{10}(1/2\zeta)$, grows without bound as $\zeta\to0$.}
\label{fig:secondbode}
\end{figure}

\begin{remark}[one fact, three pictures]
Light damping shows up three times: as poles close to the imaginary axis, as
ringing in the step response, and as a resonant peak in the Bode plot. These are
the same fact seen three ways. A pole at $-\zeta\omega_n+j\omega_d$ lies a distance
$\zeta\omega_n$ from the imaginary axis, so as $j\omega$ sweeps past it the
pole-vector of Remark~\ref{rem:vectors} becomes very short --- and a short vector
in the denominator is a large magnitude. Indeed at $\omega=\omega_d$ that vector
has length $\zeta\omega_n$ and the other pole-vector has length about $2\omega_d$,
which gives $|G|\approx\omega_n^{2}/(2\zeta\omega_n\omega_d)=M_r$.
\end{remark}

\begin{example}[resonance with no damping]
Push the limit $\zeta\to0$ all the way. With $\zeta=0$, $G(s)=\omega_n^{2}/(s^{2}+\omega_n^{2})$
has its poles on the imaginary axis, so it is not stable and \eqref{eq:freqresp}
does not apply. Drive it exactly at its natural frequency, $u=\sin\omega_n t$:
\[
   Y(s) = \frac{\omega_n^{2}}{s^{2}+\omega_n^{2}}\cdot\frac{\omega_n}{s^{2}+\omega_n^{2}}
        = \frac{\omega_n^{3}}{(s^{2}+\omega_n^{2})^{2}},
   \qquad
   y(t) = \tfrac12\bigl(\sin\omega_n t-\omega_n t\cos\omega_n t\bigr).
\]
The input's poles coincide with the plant's, the pole at $\pm j\omega_n$ becomes a
double pole, and by \S\ref{sec:pfe} a double pole brings a factor of $t$: the
oscillation grows linearly and without bound. This is resonance in its purest
form, and the infinite peak of the $\zeta\to0$ Bode curve is its frequency-domain
shadow.
\end{example}

\begin{remark}[dominant poles]\label{rem:dominant}
Why study the second-order system so closely, when real plants have many poles?
Because in a stable system each pole contributes a term $e^{p_it}$, and the terms
from poles far to the left die out long before those from poles near the imaginary
axis. The poles closest to the axis --- the \emph{dominant} poles --- therefore
shape everything one sees after the first moments. When they are a complex pair
well separated from the rest (and not nearly cancelled by a zero, which would make
their residue small), an $n$th-order system behaves very nearly like the
second-order system with those two poles. That is why $\zeta$, $\omega_n$, overshoot
and settling time serve as the working vocabulary even for plants that are not
second order.
\end{remark}

\subsection{A catalogue of elementary factors}\label{sec:catalogue}
Every rational transfer function is a product of the factors below, so every Bode
plot is a sum of their pictures. The first-order factors are written in the
normalised form that starts at $0$\,dB.

\begin{center}\small
\renewcommand{\arraystretch}{1.6}
\begin{tabular}{@{}llll@{}}
\toprule
factor & magnitude & phase & at the corner\\
\midrule
$K>0$ & flat at $20\log_{10}K$ & $0^\circ$ & ($K<0$ adds $\pm180^\circ$)\\
$s$ & $+20$\,dB/dec, $0$\,dB at $\omega=1$ & $+90^\circ$ & ---\\
$1/s$ & $-20$\,dB/dec, $0$\,dB at $\omega=1$ & $-90^\circ$ & ---\\
$1+s/z$ & $0$, then $+20$\,dB/dec from $\omega=z$ & $0^\circ\to+90^\circ$ & $+3$\,dB, $+45^\circ$\\
$1/(1+s/p)$ & $0$, then $-20$\,dB/dec from $\omega=p$ & $0^\circ\to-90^\circ$ & $-3$\,dB, $-45^\circ$\\
$\dfrac{\omega_n^{2}}{s^{2}+2\zeta\omega_n s+\omega_n^{2}}$ & $0$, then $-40$\,dB/dec from $\omega_n$ & $0^\circ\to-180^\circ$ & $1/(2\zeta)$, $-90^\circ$\\
\bottomrule
\end{tabular}
\end{center}

\begin{example}[a complete sketch]
Take $P(s)=\dfrac{10(s+1)}{s(s+10)}=\dfrac{1+s}{s\,(1+s/10)}$. It is a product of
an integrator, a first-order zero with corner $1$, and a first-order pole with corner
$10$. Reading left to right: below $\omega=1$ only the integrator is active, so the
magnitude falls at $-20$\,dB/decade through $0$\,dB at $\omega=1$ and the phase is
$-90^\circ$; at $\omega=1$ the zero cancels the slope and the magnitude flattens at
$0$\,dB while the phase rises toward $0^\circ$; at $\omega=10$ the pole brings the
slope back to $-20$\,dB/decade and the phase back toward $-90^\circ$. The whole
plot is three straight pieces joined at the corners $1$ and $10$.
\end{example}

%%%%%%%%%%%%%%%%%%%%%%%%%%%%%%%%%%%%%%%%%%%%%%%%%%%%%%%%%%%%%%%%%%%%%%%%%%%%%%
\section{The Nyquist Plot and Closed-Loop Stability}\label{sec:nyquist}

\subsection{The question, restated}
By \S\ref{sec:blocks}, the closed loop with loop transfer function $L=PC$ is stable
exactly when the characteristic equation $1+L(s)=0$ has no roots with
$\operatorname{Re}s\ge0$. Put
\[
   F(s) = 1+L(s) = \frac{D_L(s)+N_L(s)}{D_L(s)} .
\]
The \emph{zeros} of $F$ are the closed-loop poles --- unknown, and exactly what we
want to locate. The \emph{poles} of $F$ are the poles of $L$, the open-loop poles
--- known, since we know the plant and we chose the controller. The idea due to
Nyquist is to count the unknown zeros in the right half-plane without computing
them, by watching the phase of $F$ as $s$ travels around the boundary of that
half-plane.

\subsection{Counting by angles: the principle of the argument}\label{sec:argument}
Let $\Gamma$ be a simple closed curve in the $s$-plane, traversed once
\emph{clockwise}, and let $F$ be rational with no poles or zeros on $\Gamma$ itself.
As $s$ goes around $\Gamma$, the point $F(s)$ traces a closed curve $F(\Gamma)$ in
the $F$-plane (Figure~\ref{fig:argument}). Factor $F$ as in \eqref{eq:factored};
by \S\ref{sec:phase},
\[
   \angle F(s) = \angle k + \sum_i \angle(s-z_i) - \sum_j \angle(s-p_j) .
\]
Now follow a single factor. If $z_i$ lies \emph{inside} $\Gamma$, the vector
$s-z_i$ from $z_i$ to the moving point $s$ turns once all the way round as $s$
circles $\Gamma$ --- clockwise, so its angle changes by $-2\pi$. If $z_i$ lies
\emph{outside}, the vector swings to and fro but comes back to where it started:
net change $0$. The same holds for each pole, with the opposite sign in
$\angle F$. Adding up the factors:

\begin{theorem}[principle of the argument]\label{thm:argument}
Let $Z$ and $P$ be the numbers of zeros and poles of $F$ inside $\Gamma$, counted
with multiplicity. As $s$ traverses $\Gamma$ once clockwise, $\angle F(s)$ changes
by $-2\pi(Z-P)$ in total. That is, the curve $F(\Gamma)$ encircles the origin
\[
   N = Z - P
\]
times clockwise, a counterclockwise encirclement counting as $-1$.
\end{theorem}

For instance, put two points $a$, $b$ inside $\Gamma$ and a third point $c$
outside. Then $F=s-a$ winds once clockwise around the origin; $F=(s-a)(s-b)$
winds twice, because $\theta_a$ and $\theta_b$ each turn through a full circle; and
$F=(s-a)(s-b)(s-c)$ \emph{still} winds twice, because $\theta_c$ only rocks within
the wedge between the two tangents from $c$ to $\Gamma$. A pole reverses the sense:
$F=1/(s-a)$ winds once \emph{counterclockwise}, since $\angle F=-\theta_a$.

The angle count \emph{is} the proof: it is the phase rule of \S\ref{sec:phase}
applied to a point that goes around a loop instead of up the imaginary axis. (In
complex analysis the same count appears as
$\frac{1}{2\pi j}\oint_\Gamma F'(s)/F(s)\,\d s = Z-P$ for a counterclockwise
$\Gamma$: $F'/F=\sum_i 1/(s-z_i)-\sum_j 1/(s-p_j)$, and each term enclosed by
$\Gamma$ contributes a residue of $+1$ or $-1$.)

\begin{figure}[htb]
\centering
\begin{tikzpicture}[>=Stealth,
   polex/.pic={\draw[thick] (-0.11,-0.11)--(0.11,0.11) (-0.11,0.11)--(0.11,-0.11);}]
  % ---- s-plane ----
  \begin{scope}
    \node[lbl] at (-1.9,1.55) {$s$-plane};
    \draw[thick] (0,0) ellipse (1.6 and 1.1);
    \draw[->,thick] (100:1.6 and 1.1) arc[start angle=100,end angle=80,x radius=1.6,y radius=1.1];
    \draw[->,thick] (-80:1.6 and 1.1) arc[start angle=-80,end angle=-100,x radius=1.6,y radius=1.1];
    \node[lbl] at (1.35,-1.05) {$\Gamma$};
    \draw[thick] (-0.35,0.15) circle (0.1);
    \node[lbl,below] at (-0.35,0.03) {$z_1$};
    \pic at (2.35,0.75) {polex};
    \node[lbl,right] at (2.45,0.75) {$p_1$};
    \coordinate (s) at (40:1.6 and 1.1);
    \fill (s) circle (1.3pt); \node[lbl,above right] at (s) {$s$};
    \draw[->] (-0.35,0.15) -- ($(s)+(-0.05,-0.03)$);
    \draw[->,gray] (2.35,0.75) -- ($(s)+(0.06,0.0)$);
  \end{scope}
  % ---- F-plane ----
  \begin{scope}[xshift=6.4cm]
    \node[lbl] at (-1.6,1.55) {$F$-plane};
    \draw[->] (-1.9,0) -- (1.9,0);
    \draw[->] (0,-1.5) -- (0,1.6);
    \draw[thick] plot[domain=0:360,samples=90,smooth,variable=\a]
      ({0.25+1.25*cos(\a)+0.25*cos(2*\a)},{-0.95*sin(\a)+0.12*sin(3*\a)});
    \draw[->,thick] (1.5,-0.1) -- (1.48,-0.3);
    \fill (0,0) circle (1.3pt);
    \node[lbl] at (1.45,1.05) {$F(\Gamma)$};
  \end{scope}
\end{tikzpicture}
\caption{The principle of the argument. As $s$ goes once clockwise around $\Gamma$,
the vector $s-z_1$ from the enclosed zero turns once clockwise, while the vector
$s-p_1$ from the pole outside only rocks back and forth. With one zero and no pole
inside, the image $F(\Gamma)$ encircles the origin once clockwise: $N=Z-P=1$.}
\label{fig:argument}
\end{figure}

\subsection{The Nyquist contour and the Nyquist plot}
To count right-half-plane roots, take $\Gamma$ to be the boundary of the right
half-plane. The \emph{Nyquist contour} $\Gamma_N$ runs up the imaginary axis from
$-jR$ to $+jR$ and returns along the semicircle $|s|=R$ through the right
half-plane, clockwise, with $R\to\infty$ (Figure~\ref{fig:contour}). Its inside is
the whole open right half-plane. If $L$ has poles \emph{on} the imaginary axis ---
an integrator $1/s$, say --- $\Gamma_N$ is not allowed to pass through them. It
detours around each one on a small semicircle of radius $\varepsilon\to0$ that
bulges into the right half-plane, which leaves those poles outside.

\begin{figure}[htb]
\centering
\begin{tikzpicture}[>=Stealth,scale=0.95,
   polex/.pic={\draw[thick] (-0.11,-0.11)--(0.11,0.11) (-0.11,0.11)--(0.11,-0.11);}]
  \fill[black!5] (0,-2.2) -- (0,-0.32) arc (-90:90:0.32) -- (0,2.2) arc (90:-90:2.2) -- cycle;
  \draw[->] (-1.4,0) -- (3.0,0) node[right] {$\operatorname{Re}s$};
  \draw[->] (0,-2.7) -- (0,2.75) node[above] {$\operatorname{Im}s$};
  \draw[thick] (0,-2.2) -- (0,-0.32);
  \draw[->,thick] (0,-2.2) -- (0,-1.15);
  \draw[thick] (0,-0.32) arc (-90:90:0.32);
  \draw[thick] (0,0.32) -- (0,2.2);
  \draw[->,thick] (0,0.32) -- (0,1.35);
  \draw[thick] (0,2.2) arc (90:-90:2.2);
  \draw[->,thick] (45:2.2) arc (45:35:2.2);
  \pic at (0,0) {polex};
  \node[lbl,anchor=west] at (1.6,1.95) {$R\to\infty$};
  \node[lbl,anchor=west] at (0.42,0.36) {$\varepsilon\to0$};
  \node[lbl] at (1.05,-0.95) {RHP};
\end{tikzpicture}
\caption{The Nyquist contour $\Gamma_N$, traversed clockwise: up the imaginary axis,
then back along a semicircle of infinite radius. Its inside (shaded) is the open
right half-plane. A pole on the imaginary axis --- here at the origin --- is
avoided by a small detour into the right half-plane.}
\label{fig:contour}
\end{figure}

The image of $\Gamma_N$ under $L$ is the \emph{Nyquist plot} of $L$. Two facts make
it easy to draw.
\begin{itemize}[leftmargin=*,itemsep=2pt]
\item If $L$ is strictly proper, $L(s)\to0$ as $|s|\to\infty$, so the whole large
semicircle maps to the single point $0$. The plot is then just the curve
$L(j\omega)$, $-\infty<\omega<\infty$ --- the frequency response, drawn in the
complex plane instead of as two Bode graphs.
\item Because $L$ has real coefficients, $L(-j\omega)=\overline{L(j\omega)}$: the
half for $\omega<0$ is the mirror image, across the real axis, of the half for
$\omega>0$.
\end{itemize}
In practice one plots $L(j\omega)$ for $\omega\ge0$ --- read straight off the Bode
plot, since the distance from the origin is $|L(j\omega)|$ and the angle is
$\angle L(j\omega)$ --- reflects it in the real axis, and adds whatever the small
detours contribute.

\subsection{The Nyquist stability criterion}
Apply Theorem~\ref{thm:argument} to $F=1+L$ on $\Gamma_N$. The zeros of $F$ inside
$\Gamma_N$ are the closed-loop poles in the right half-plane; the poles of $F$ inside
are the open-loop poles there. And $F(\Gamma_N)$ is the Nyquist plot of $L$ shifted
one unit to the right, so encircling the origin with $F$ is the same as encircling
the point $-1$ with $L$. That is the whole criterion.

\begin{keybox}
\textbf{Nyquist stability criterion.} Let $P$ be the number of poles of $L(s)$ in
the open right half-plane, and let $N$ be the number of times the Nyquist plot of
$L$ encircles the point $-1$ clockwise (counterclockwise encirclements counted
negative). Then the number of closed-loop poles in the right half-plane is
\[
   Z = N + P .
\]
The closed loop is stable if and only if $Z=0$, that is $N=-P$, and the plot does
not pass through $-1$ (if it does, the closed loop has a pole on the imaginary
axis). In particular, if $L$ is itself stable ($P=0$), the closed loop is stable if
and only if the Nyquist plot does not encircle $-1$.

\smallskip\noindent
Counted the other way round, the same statement reads: \emph{the number of
counterclockwise encirclements of $-1$ equals $P-Z$.}
\end{keybox}

\begin{remark}[why this is worth the trouble]
Two things set the criterion apart from simply computing the roots of $D_L+N_L$.
First, it needs only $L(j\omega)$, and that can be \emph{measured}: drive the open
loop with sinusoids and record gain and phase, with no model of the plant at all.
Second, it shows not only whether the loop is stable but how \emph{close} it is to
instability. A plot that passes near $-1$ can be pushed across it by a small change
of gain or phase --- by exactly the kind of modelling error that \S1 said we must
expect.
\end{remark}

\subsection{Examples}
\begin{example}[a stable open loop]\label{ex:nyq1}
$L(s)=\dfrac{K}{(s+1)(s+2)(s+3)}$ with $K>0$. All open-loop poles are in the left
half-plane, so $P=0$. At $s=j\omega$ the denominator is
\[
   (6-6\omega^{2}) + j\,(11\omega-\omega^{3}) .
\]
The plot starts at $L(0)=K/6$ on the positive real axis and moves clockwise through
the fourth and third quadrants as the phase falls from $0^\circ$ toward
$-270^\circ$. It crosses the negative real axis where the imaginary part of the
denominator vanishes, at $\omega^{2}=11$, and there $L=K/(6-66)=-K/60$
(Figure~\ref{fig:nyqex}, left, drawn for $K=20$). If $K<60$ the crossing lies to the
right of $-1$ and there is no encirclement: $N=0$, $Z=0$, stable. If $K>60$ the plot
encircles $-1$ twice clockwise: $N=2$, so two closed-loop poles lie in the right
half-plane. At the boundary $K=60$ the plot passes through $-1$ at
$\omega=\sqrt{11}$, and indeed
\[
   D_L+N_L = s^{3}+6s^{2}+11s+66 = (s+6)(s^{2}+11)
\]
has closed-loop poles exactly at $\pm j\sqrt{11}$ --- the frequency at which the plot
meets $-1$.
\end{example}

\begin{example}[an integrator in the loop]\label{ex:nyq2}
$L(s)=\dfrac{K}{s(s+1)}$ with $K>0$. The pole at the origin lies on the imaginary
axis, so $\Gamma_N$ detours around it on $s=\varepsilon e^{j\phi}$ with $\phi$
running from $-90^\circ$ to $+90^\circ$. There $L(s)\approx K/s =
(K/\varepsilon)e^{-j\phi}$: a semicircle of enormous radius $K/\varepsilon$, swept
from $+90^\circ$ to $-90^\circ$, that is clockwise through the right half of the
plane. On the imaginary axis,
\[
   L(j\omega) = -\frac{K}{1+\omega^{2}} - j\,\frac{K}{\omega(1+\omega^{2})},
\]
so for $\omega>0$ the plot rises from far below along the asymptote
$\operatorname{Re}L=-K$, stays in the third quadrant, and enters the origin with
phase $-180^\circ$ as $\omega\to\infty$. Its imaginary part never vanishes, so it
never crosses the negative real axis and $-1$ is never encircled. The detoured pole
is not counted, so $P=0$; $N=0$; hence $Z=0$ and the loop is stable for every
$K>0$. Directly: $D_L+N_L=s^{2}+s+K$, whose roots have negative real part for all
$K>0$.
\end{example}

The detour always behaves as in this example. Near a simple pole $p$ on the
imaginary axis, $L(s)\approx c/(s-p)$, and as the detour turns $s-p$ through
$+180^\circ$ (from pointing down to pointing up, by way of the right), the image
turns through $-180^\circ$ at infinite radius: half a clockwise circle. A double
pole on the axis gives a full clockwise circle, $-360^\circ$ --- enough, by itself,
to wrap around $-1$.

\begin{example}[a pair of poles on the imaginary axis]\label{ex:nyq4}
$L(s)=\dfrac{1}{s^{2}+4}=\dfrac{1}{(s-2j)(s+2j)}$. Both poles lie on the axis and are
detoured, so $P=0$. Away from them, $L(j\omega)=1/(4-\omega^{2})$ is real: it runs
along the positive real axis from $\tfrac14$ out to $+\infty$ as $|\omega|\to2$, and
along the negative real axis from $-\infty$ back to $0$ for $|\omega|>2$. Each detour
contributes a clockwise half-circle at infinity --- the one at $+2j$ through the
lower half-plane, the one at $-2j$ through the upper --- and together they close the
plot into an infinitely large circle. The whole negative real axis is covered, so
the plot passes \emph{through} $-1$, at $4-\omega^{2}=-1$, i.e.\ $\omega=\sqrt5$. By
the criterion the closed loop has a pole on the imaginary axis, and indeed
$1+L=0$ means $s^{2}+5=0$, $s=\pm j\sqrt5$: a sustained oscillation at
$\sqrt5$\,rad/s, neither decaying nor growing.
\end{example}

\begin{example}[an unstable open loop]\label{ex:nyq3}
$L(s)=\dfrac{K}{s-1}$ with $K>0$. Now the open loop has a pole at $s=1$, so $P=1$.
On the imaginary axis
\[
   L(j\omega) = \frac{K}{j\omega-1}
   = -\frac{K}{1+\omega^{2}} - j\,\frac{K\omega}{1+\omega^{2}},
\]
which traces the circle of radius $K/2$ centred at $-K/2$: from the origin
($\omega=-\infty$) through $-\tfrac K2+j\tfrac K2$, then $-K$ ($\omega=0$), then
$-\tfrac K2-j\tfrac K2$, and back to the origin --- counterclockwise
(Figure~\ref{fig:nyqex}, right). If $K>1$ the circle contains $-1$ and goes round it
once counterclockwise: $N=-1$, so $Z=N+P=0$ and the loop is stable. If $K<1$ it
misses $-1$: $N=0$, $Z=1$, unstable. Directly, $D_L+N_L=s-1+K$, stable exactly when
$K>1$. An unstable plant \emph{needs} an encirclement to be stabilised; ``the plot
must not encircle $-1$'' is the rule only for stable open loops.
\end{example}

\begin{remark}[reading the criterion off a Bode plot]\label{rem:bodenyq}
When $L$ is stable ($P=0$) and its phase falls through $-180^\circ$ just once, at a
frequency $\omega_{pc}$, the Nyquist plot for $\omega>0$ crosses the negative real
axis just once, at the point $-|L(j\omega_{pc})|$. It encircles $-1$ exactly when
that point lies to the left of $-1$. So the criterion becomes a statement one can
check on the Bode plot:
\[
   \text{closed loop stable}
   \quad\Longleftrightarrow\quad
   |L(j\omega_{pc})| < 1
   \quad\text{where } \angle L(j\omega_{pc}) = -180^\circ .
\]
By the time the phase has turned the loop signal fully upside down, the gain must
already be below one. In Example~\ref{ex:nyq1} the phase reaches $-180^\circ$ at
$\omega_{pc}=\sqrt{11}$, where $|L|=K/60$ --- the same condition $K<60$ again. If
the gain is too high there, it can be reduced (the idea behind a lag compensator),
or the phase can be pushed back up so that $-180^\circ$ is reached later (a lead
compensator).
\end{remark}

\begin{figure}[htb]
\centering
\begin{tikzpicture}[>=Stealth,
   declare function={Lm(\w)=20/sqrt(\w*\w+1)/sqrt(\w*\w+4)/sqrt(\w*\w+9);
                     La(\w)=-(atan(\w)+atan(\w/2)+atan(\w/3));}]
  % ---------------- Example: K/((s+1)(s+2)(s+3)), K=20 ----------------
  \begin{scope}[scale=0.8]
    \draw[->] (-1.7,0) -- (4.2,0) node[right,lbl] {$\operatorname{Re}L$};
    \draw[->] (0,-2.8) -- (0,2.8) node[above,lbl] {$\operatorname{Im}L$};
    \draw[thick] plot[domain=-2:1.6,samples=120,smooth]
      ({Lm(pow(10,\x))*cos(La(pow(10,\x)))},{Lm(pow(10,\x))*sin(La(pow(10,\x)))});
    \draw[dashed] plot[domain=-2:1.6,samples=120,smooth]
      ({Lm(pow(10,\x))*cos(La(pow(10,\x)))},{-Lm(pow(10,\x))*sin(La(pow(10,\x)))});
    \draw[->,thick] ({Lm(0.30)*cos(La(0.30))},{Lm(0.30)*sin(La(0.30))})
                 -- ({Lm(0.34)*cos(La(0.34))},{Lm(0.34)*sin(La(0.34))});
    \fill (-1,0) circle (2.2pt);
    \node[lbl,above left] at (-1,0.05) {$-1$};
    \draw (-0.333,0.1) -- (-0.333,-0.1);
    \node[lbl,below left] at (-0.30,-0.08) {$-\tfrac13$};
    \node[lbl,above right] at (3.333,0) {$\tfrac{10}{3}$};
    \node[lbl,right] at (0.12,-2.05) {$\omega=1$};
    \node[lbl] at (3.35,-1.45) {$\omega>0$};
    \node[lbl] at (3.35,1.45) {$\omega<0$};
  \end{scope}
  % ---------------- Example: K/(s-1), K=1.5 ----------------
  \begin{scope}[xshift=8.6cm,scale=1.45]
    \draw[->] (-2.0,0) -- (0.7,0) node[right,lbl] {$\operatorname{Re}L$};
    \draw[->] (0,-1.1) -- (0,1.15) node[above,lbl] {$\operatorname{Im}L$};
    \draw[thick] (-0.75,0) circle (0.75);
    \draw[->,thick] (-0.75,0.75) arc (90:110:0.75);
    \draw[->,thick] (-0.75,-0.75) arc (270:290:0.75);
    \fill (-1,0) circle (1.5pt);
    \node[lbl,above] at (-1.02,0.04) {$-1$};
    \node[lbl,below left] at (-1.5,0) {$-K$};
    \node[lbl,left] at (-1.52,0.18) {$\omega=0$};
    \node[lbl,above right] at (0,0) {$\omega\to\pm\infty$};
  \end{scope}
\end{tikzpicture}
\caption{Left: the Nyquist plot of Example~\ref{ex:nyq1} for $K=20$ (solid $\omega>0$,
dashed $\omega<0$). It crosses the negative real axis at $-K/60=-\tfrac13$, to the
right of $-1$, so $-1$ is not encircled and the closed loop is stable. Right:
Example~\ref{ex:nyq3} for $K=1.5$. The open loop is unstable ($P=1$), and the plot
goes once counterclockwise around $-1$ ($N=-1$), which is exactly what makes the
closed loop stable.}
\label{fig:nyqex}
\end{figure}

%% ==== Part V ===============================================================
\partdivider{Part V --- Linear Algebra for State Space}

\lecturebreak{Lectures 7--10}{23 September -- 7 October 2026}

\section{Vector Spaces, Rank, and Linear Equations}\label{sec:vspace}

A state is a point in $\R^{n}$, an input pushes it along certain directions, and an
output reports certain combinations of its coordinates. Everything the rest of the
course does with $\dot x=Ax+Bu$ is therefore a statement about vectors, subspaces
and matrices. This Part collects the linear algebra that will be used, in the
order in which it is needed.

\subsection{Matrices, read by columns}
A matrix $A\in\R^{m\times n}$ has $m$ rows and $n$ columns; it is \emph{square} if
$m=n$, \emph{tall} if $m>n$ and \emph{fat} if $m<n$. Its transpose $A\tp\in\R^{n\times m}$
satisfies $(ABC)\tp=C\tp B\tp A\tp$, and for complex matrices the
\emph{Hermitian} (conjugate) transpose is $A^{H}=\overline{A}\tp$. The product
$C=AB$ of $A\in\R^{m\times n}$ and $B\in\R^{n\times k}$ has entries
$c_{ij}=\sum_{l}a_{il}b_{lj}$, but three other ways of reading it are far more
useful. Write $a_1,\dots,a_n$ for the columns of $A$ and $b_1\tp,\dots,b_n\tp$ for
the rows of $B$.
\begin{enumerate}[label=(\roman*),itemsep=2pt]
\item $Ax = x_1a_1+\dots+x_na_n$: a matrix times a vector is a
\emph{linear combination of the columns}.
\item $y\tp B = y_1b_1\tp+\dots+y_nb_n\tp$: a row vector times a matrix is a linear
combination of the rows.
\item $AB = \sum_{i=1}^{n} a_i\,b_i\tp$: a product is a \emph{sum of outer products},
each $a_ib_i\tp$ a matrix of rank one.
\end{enumerate}
Reading (i) is the one to keep in mind throughout: asking whether $Ax=b$ can be
solved is asking whether $b$ is a combination of the columns of $A$.

\subsection{Vector spaces, span, basis, dimension}
\begin{definition}[vector space]
A set $S$ of vectors is a \emph{vector space} (a \emph{subspace} of $\R^{n}$) if
it is closed under addition and scalar multiplication: $x,y\in S$ and
$a,b\in\R$ imply $ax+by\in S$.
\end{definition}
Taking $a=b=0$ shows that every vector space contains the origin. A line or a plane
\emph{through the origin} is a subspace; one that misses the origin is not.

A \emph{linear combination} of $a_1,\dots,a_k$ is $\sum_i c_ia_i$, and the vectors
are \emph{linearly independent} if $\sum_ic_ia_i=0$ forces every $c_i=0$ --- no one
of them can be built from the others; they all point in genuinely different
directions. The set of all linear combinations is the \emph{span},
\[
   \operatorname{span}\{a_1,\dots,a_k\} = \Bigl\{\,\sum_{i=1}^{k}c_ia_i \;:\; c_i\in\R\Bigr\},
\]
and it is a vector space: adding two combinations, or scaling one, gives another
combination. A \emph{basis} of a vector space is a set of linearly independent
vectors that spans it. A basis is not unique, but the number of vectors in every
basis is the same; that number is the \emph{dimension}. In the plane $x_3=0$ of
$\R^{3}$, a single vector cannot span the plane, while any two non-parallel vectors
in it form a basis, so the plane has dimension $2$. A basis whose vectors are
mutually perpendicular and of unit length is \emph{orthonormal}.

There are two ways to describe a subspace, and both will be used.
\begin{itemize}[leftmargin=*,itemsep=2pt]
\item \textbf{By building it up} as a span. The plane $x_3=0$ is
$\operatorname{span}\{e_1,e_2\}$; each independent generator adds one to the
dimension ($0\to2$).
\item \textbf{By cutting it down} with equations. The same plane is
$\{x\in\R^{3}:[0\;0\;1]\,x=0\}$; each independent constraint removes one dimension
($3\to2$).
\end{itemize}
The line through the origin at $45^\circ$ in the $x_2x_3$-plane is
$\operatorname{span}\{(0,1,1)\}$, or equally
$\{x\in\R^{3}: [1\;0\;0]\,x=0,\ [0\;1\;-1]\,x=0\}$: two independent constraints
leave $3-2=1$ dimension.

\subsection{Range, null space, and the equation \texorpdfstring{$Ax=b$}{Ax=b}}
\label{sec:rangenull}
Two subspaces belong to every matrix $A\in\R^{m\times n}$, one of each kind:
\begin{equation}\label{eq:rangenull}
   \mathcal R(A) = \{Ax : x\in\R^{n}\}\subseteq\R^{m},
   \qquad
   \mathcal N(A) = \{x\in\R^{n} : Ax=0\}\subseteq\R^{n}.
\end{equation}
By reading (i), the \emph{range} $\mathcal R(A)$ is the span of the columns --- it
is also called the \emph{column space}, a description by building up. The
\emph{null space} $\mathcal N(A)$ is a description by cutting down. The
\emph{rank} of $A$ is the number of linearly independent columns, which is always
equal to the number of linearly independent rows, and it is the dimension of
$\mathcal R(A)$.

The two subspaces answer the two questions one asks about a linear equation
$Ax=b$:
\begin{itemize}[leftmargin=*,itemsep=2pt]
\item \textbf{Existence.} A solution exists if and only if $b\in\mathcal R(A)$ ---
if and only if $b$ is a combination of the columns.
\item \textbf{Uniqueness.} If $Ax=b$ and $Ax^{*}=b$, then $A(x-x^{*})=0$, so the
two solutions differ by an element of $\mathcal N(A)$. A solution is unique if
and only if $\mathcal N(A)=\{0\}$, which by reading (i) says that
$x_1a_1+\dots+x_na_n=0$ only for $x=0$: the columns are linearly independent.
\end{itemize}
A matrix whose columns are independent, $\rank A=n$, is said to have \emph{full
column rank}, and then $\mathcal N(A)=\{0\}$. Existence and uniqueness are
different questions with different answers, and it pays to keep them apart.

\begin{remark}[counting dimensions]
The cutting-down count above is a general fact. For $A\in\R^{m\times n}$,
\[
   \dim\mathcal R(A)+\dim\mathcal N(A) = n ,
\]
the \emph{rank--nullity} relation: of the $n$ input directions, $\rank A$ of them
reach independent outputs and the remaining $n-\rank A$ are sent to zero.
\end{remark}

\begin{example}[the equilibria of the two carts]\label{ex:cartnull}
For the two-cart matrix $A$ of \eqref{eq:ABCD}, the states that stay put when no
force is applied are the solutions of $Ax=0$. Rows $1$ and $3$ give $x_2=x_4=0$:
both carts at rest. Row $2$ then reads $\frac{k}{m_1}(x_3-x_1)=0$, so $x_1=x_3$, and
row $4$ says the same. Hence
\[
   \mathcal N(A) = \operatorname{span}\{(1,0,1,0)\tp\},
   \qquad \rank A = 3 .
\]
Any position with both carts equally far along and the spring unstretched is an
equilibrium --- there is no ground spring to pull them back. In particular the
target of Question~\ref{q:goal}, both carts $1$\,m ahead and at rest,
$x=(1,0,1,0)\tp$, is an equilibrium of the uncontrolled system: once there, the
carts stay there with $u=0$. Keeping the damper $b$ changes nothing, since every
$b$-term is multiplied by a velocity.
\end{example}

\subsection{Inverse and determinant}
A square matrix $A\in\R^{n\times n}$ is \emph{invertible} if there is $A^{-1}$ with
$AA^{-1}=A^{-1}A=I$. Writing $A^{-1}=[x_1\;\cdots\;x_n]$, the equation $AA^{-1}=I$ is
the $n$ equations $Ax_i=e_i$, so an inverse exists exactly when every $e_i$, and
hence all of $\R^{n}$, lies in $\mathcal R(A)$; for a square matrix that happens
exactly when $\mathcal N(A)=\{0\}$. The following are therefore equivalent:
$A$ is invertible; $\mathcal R(A)=\R^{n}$; $\mathcal N(A)=\{0\}$; the columns are
independent; $\rank A=n$; $\det A\neq0$; $0$ is not an eigenvalue of $A$.

For a $2\times2$ matrix,
\[
   \begin{bmatrix}a&b\\c&d\end{bmatrix}^{-1}
   = \frac{1}{ad-bc}\begin{bmatrix}d&-b\\-c&a\end{bmatrix}
   = \frac{\operatorname{adj}A}{\det A},
\]
and the formula $A^{-1}=\operatorname{adj}A/\det A$ holds in every dimension; it is
the formula behind $(sI-A)^{-1}=\operatorname{adj}(sI-A)/\det(sI-A)$ of
\S\ref{sec:ss2tf}. The determinant satisfies $\det A\tp=\det A$ and
$\det(ABC)=\det A\,\det B\,\det C$.

\begin{proposition}[Schur complement]\label{prop:schur}
Let $M=\begin{bmatrix}A&B\\C&D\end{bmatrix}$ with $A$ and $D$ square. If $A$ is
invertible,
\[
   \det M = \det A\,\det\bigl(D-CA^{-1}B\bigr);
\]
if $D$ is invertible, $\det M=\det D\,\det\bigl(A-BD^{-1}C\bigr)$.
\end{proposition}
\begin{proof}
Block elimination: $\begin{bmatrix}A&B\\C&D\end{bmatrix}
=\begin{bmatrix}I&0\\CA^{-1}&I\end{bmatrix}
\begin{bmatrix}A&B\\0&D-CA^{-1}B\end{bmatrix}$. The first factor has determinant
$1$ and the second is block triangular. The other case is symmetric.
\end{proof}
The matrix $D-CA^{-1}B$ is the \emph{Schur complement} of $A$ in $M$ --- what is
left of $D$ after the block $A$ has been eliminated, exactly as in Gaussian
elimination on scalars.

%%%%%%%%%%%%%%%%%%%%%%%%%%%%%%%%%%%%%%%%%%%%%%%%%%%%%%%%%%%%%%%%%%%%%%%%%%%%%%
\section{Coordinates, Eigenvalues, and Diagonalization}\label{sec:eigen}

\subsection{Changing state coordinates}\label{sec:similarity}
The state vector of a system is a choice. Positions and velocities of each cart
are one choice; the position of the pair and the stretch of the spring would be
another. Any invertible $T\in\R^{n\times n}$ defines new coordinates
\[
   z(t) = T\,x(t) .
\]
Substituting $x=T^{-1}z$ into \eqref{eq:std},
\begin{equation}\label{eq:simtransform}
   \dot z = TAT^{-1}\,z + TB\,u, \qquad y = CT^{-1}\,z + D\,u .
\end{equation}
The quadruple $(A,B,C,D)$ becomes $(TAT^{-1},TB,CT^{-1},D)$, a
\emph{similarity transformation}. The physics has not changed, so neither should
anything one can measure --- and indeed the transfer function is unchanged:
\[
   CT^{-1}\bigl(sI-TAT^{-1}\bigr)^{-1}TB + D
   = CT^{-1}\bigl(T(sI-A)T^{-1}\bigr)^{-1}TB + D
   = C(sI-A)^{-1}B + D ,
\]
using $sI=T(sI)T^{-1}$. This is one reason the realization of
\S\ref{sec:ss2tf} is far from unique: every invertible $T$ turns one realization
into another with the same transfer function.

\subsection{Eigenvalues and eigenvectors}
\begin{definition}[eigenvalue, eigenvector]
A number $\lambda$ is an \emph{eigenvalue} of $A\in\R^{n\times n}$, with
\emph{eigenvector} $v\neq0$, if $Av=\lambda v$.
\end{definition}
An eigenvector is a direction that $A$ does not turn, only stretches by the factor
$\lambda$. Rewriting $Av=\lambda v$ as $(\lambda I-A)v=0$ shows that $\lambda$ is an
eigenvalue exactly when $\mathcal N(\lambda I-A)\neq\{0\}$, i.e.\ when $\lambda I-A$
is not invertible:
\begin{equation}\label{eq:charpoly}
   \det(\lambda I-A) = \lambda^{n}+a_{n-1}\lambda^{n-1}+\dots+a_1\lambda+a_0 = 0 .
\end{equation}
The left side is the \emph{characteristic polynomial}, of degree $n$, so $A$ has
$n$ eigenvalues counted with multiplicity --- complex ones in conjugate pairs,
since the coefficients are real. Every nonzero vector in
$\mathcal N(\lambda I-A)$ is an eigenvector, and that subspace is the
\emph{eigenspace} of $\lambda$.

Eigenvalues do not depend on the coordinates: by \eqref{eq:simtransform},
\[
   \det\bigl(\lambda I-TAT^{-1}\bigr) = \det\bigl(T(\lambda I-A)T^{-1}\bigr)
   = \det T\,\det(\lambda I-A)\,\det T^{-1} = \det(\lambda I-A),
\]
so the whole characteristic polynomial is invariant. And the denominator of
$P(s)=C\operatorname{adj}(sI-A)B/\det(sI-A)+D$ is exactly $\det(sI-A)$: \emph{the
eigenvalues of $A$ are the poles of the transfer function}, except where a factor
cancels against the numerator. The two descriptions of Figure~\ref{fig:routes}
meet here.

\subsection{Diagonalization}\label{sec:diag}
Suppose $A$ has $n$ linearly independent eigenvectors $v_1,\dots,v_n$ with
eigenvalues $\lambda_1,\dots,\lambda_n$. Stacking the $n$ equations
$Av_i=\lambda_iv_i$ side by side,
\[
   A\,[v_1\;\cdots\;v_n] = [\lambda_1v_1\;\cdots\;\lambda_nv_n]
   = [v_1\;\cdots\;v_n]
     \begin{bmatrix}\lambda_1&&\\&\ddots&\\&&\lambda_n\end{bmatrix},
\]
that is
\begin{equation}\label{eq:diagonalize}
   AV = V\Lambda,\qquad V^{-1}AV=\Lambda,\qquad A = V\Lambda V^{-1}.
\end{equation}
$V$ is invertible precisely because its columns are independent. In the language of
\S\ref{sec:similarity}, the coordinates $z=V^{-1}x$ make the state matrix diagonal.

When are there $n$ independent eigenvectors? Two counts decide it. The
\emph{algebraic multiplicity} of an eigenvalue is how many times it is a root of
\eqref{eq:charpoly}; its \emph{geometric multiplicity} is
$\dim\mathcal N(\lambda I-A)$, the number of independent eigenvectors it has.
Always $1\le$ geometric $\le$ algebraic.
\begin{itemize}[leftmargin=*,itemsep=2pt]
\item Eigenvectors belonging to \emph{distinct} eigenvalues are linearly
independent. So if all $n$ eigenvalues are distinct, $A$ is diagonalizable.
\item In general, $A$ is diagonalizable if and only if the geometric multiplicity
equals the algebraic multiplicity for every eigenvalue.
\end{itemize}

\begin{example}[equal multiplicities, and not]
Both
\[
   A=\begin{bmatrix}1&0&0&0\\0&2&0&0\\0&0&2&0\\0&0&0&3\end{bmatrix}
   \quad\text{and}\quad
   B=\begin{bmatrix}1&0&0&0\\0&2&1&0\\0&0&2&0\\0&0&0&3\end{bmatrix}
\]
have eigenvalues $1,2,2,3$; the eigenvalue $2$ has algebraic multiplicity $2$ in
each. For $A$, $\mathcal N(2I-A)=\operatorname{span}\{e_2,e_3\}$ has dimension $2$:
two independent eigenvectors, and $A$ is (already) diagonal. For $B$,
\[
   2I-B=\begin{bmatrix}1&0&0&0\\0&0&-1&0\\0&0&0&0\\0&0&0&-1\end{bmatrix},
   \qquad \mathcal N(2I-B)=\operatorname{span}\{e_2\},
\]
of dimension $1$: the two copies of the eigenvalue share a single eigenvector
direction, $V$ cannot be made invertible, and $B$ is not diagonalizable. The single
off-diagonal $1$ is what makes the difference.
\end{example}

\begin{example}[the two carts, decoupled]\label{ex:cartdecouple}
Take equal masses $m_1=m_2=m$ (and $b=0$, as in \eqref{eq:ABCD}). Adding and
subtracting the two equations of motion \eqref{eq:eom},
\[
   m(\ddot p_1+\ddot p_2) = u - c(\dot p_1+\dot p_2),\qquad
   m(\ddot p_1-\ddot p_2) = u - 2k(p_1-p_2) - c(\dot p_1-\dot p_2).
\]
So the coordinates $z=(p_1+p_2,\;\dot p_1+\dot p_2,\;p_1-p_2,\;\dot p_1-\dot p_2)$,
i.e.\ $z=Tx$ with
\[
   T=\begin{bmatrix}1&0&1&0\\0&1&0&1\\1&0&-1&0\\0&1&0&-1\end{bmatrix},
   \qquad
   TAT^{-1}=\begin{bmatrix}0&1&0&0\\0&-\frac cm&0&0\\0&0&0&1\\0&0&-\frac{2k}m&-\frac cm\end{bmatrix},
   \qquad TB=\frac1m\begin{bmatrix}0\\1\\0\\1\end{bmatrix},
\]
split the system into two independent pieces: the motion of the pair as a whole,
which feels only friction, and the oscillation of the carts against each other,
which feels friction and twice the spring. The eigenvalues can now be read off
block by block: $0$ and $-c/m$ from the first, and the roots of
$\lambda^{2}+\frac cm\lambda+\frac{2k}{m}$ from the second. For the parameters used
in the first problem set ($m=1$, $k=4$, $c=0.4$) they are $0$, $-0.4$ and
$-0.2\pm2.8213j$. The eigenvalue $0$ belongs to the eigenvector $(1,0,1,0)\tp$ of
Example~\ref{ex:cartnull}: the free drift of the pair, which nothing in the
open loop holds back.
\end{example}

\subsection{Symmetric matrices}\label{sec:symmetric}
Symmetric matrices, $A=A\tp$, are the best behaved of all, and they will return in
\S\ref{sec:svd}.
\begin{theorem}[spectral theorem]\label{thm:spectral}
Let $A=A\tp\in\R^{n\times n}$. Then every eigenvalue of $A$ is real, eigenvectors
belonging to distinct eigenvalues are orthogonal, and one can choose $n$
orthonormal eigenvectors. With them as the columns of $V$,
\[
   V\tp V=I,\qquad V^{-1}=V\tp,\qquad
   A = V\Lambda V\tp = \sum_{i=1}^{n}\lambda_i\,v_iv_i\tp .
\]
\end{theorem}
\begin{proof}[Orthogonality]
If $Av_i=\lambda_iv_i$ and $Av_j=\lambda_jv_j$, then
$\lambda_i\,v_j\tp v_i = v_j\tp Av_i = v_j\tp A\tp v_i = (Av_j)\tp v_i
= \lambda_j\,v_j\tp v_i$, so $(\lambda_i-\lambda_j)\,v_j\tp v_i=0$, and
$v_j\tp v_i=0$ when $\lambda_i\neq\lambda_j$.
\end{proof}
Within an eigenspace of dimension $k$ one can always pick $k$ orthonormal vectors,
and eigenspaces of different eigenvalues are orthogonal to one another, so all the
choices together give an orthonormal set. For a symmetric matrix the inverse of
$V$ costs nothing --- it is the transpose --- and $V\tp AV=\Lambda$ is much easier
to compute than $V^{-1}AV$. The last form, $A=\sum\lambda_iv_iv_i\tp$, writes $A$ as a
sum of rank-one pieces, one along each eigen-direction.

\subsection{The Jordan form}\label{sec:jordan}
A matrix with too few eigenvectors, such as $B$ above, cannot be diagonalized, but
every square matrix can be brought to the next best thing.
\begin{theorem}[Jordan form]
Every $A\in\R^{n\times n}$ satisfies $AV=VJ$ for some invertible (possibly complex)
$V$, where $J$ is block diagonal with \emph{Jordan blocks}
\[
   J_k(\lambda)=\begin{bmatrix}\lambda&1&&\\&\lambda&\ddots&\\&&\ddots&1\\&&&\lambda\end{bmatrix}
\]
on the diagonal. $A$ is diagonalizable exactly when every block is $1\times1$.
\end{theorem}
To see what the columns of $V$ are, take a $3\times3$ example,
\[
   AV=VJ,\qquad J=\begin{bmatrix}\lambda_1&0&0\\0&\lambda_2&1\\0&0&\lambda_2\end{bmatrix},
   \qquad
   [Av_1\;Av_2\;Av_3]=[\lambda_1v_1\;\;\lambda_2v_2\;\;\lambda_2v_3+v_2].
\]
The first two columns are ordinary eigenvectors. The third satisfies
\[
   (A-\lambda_2I)\,v_3 = v_2 ,
\]
so it is not an eigenvector --- $A-\lambda_2I$ does not annihilate it, but sends it
onto the eigenvector $v_2$. Such a $v_3$ is a \emph{generalized eigenvector}. A
generalized eigenvector exists exactly when some eigenvalue has algebraic
multiplicity greater than its geometric multiplicity; it supplies the direction
that the missing eigenvector would have provided.

%%%%%%%%%%%%%%%%%%%%%%%%%%%%%%%%%%%%%%%%%%%%%%%%%%%%%%%%%%%%%%%%%%%%%%%%%%%%%%
\section{The Matrix Exponential and the State Response}\label{sec:expm}

\subsection{Solving \texorpdfstring{$\dot x=Ax$}{xdot = Ax}}
For a scalar, $\dot x=ax$ has the solution $x(t)=e^{at}x_0$. The same is true for
a system, once the exponential of a matrix is defined by the same series.
\begin{definition}[matrix exponential]
For $A\in\R^{n\times n}$,
\[
   e^{At} = I + At + \frac{(At)^{2}}{2!} + \frac{(At)^{3}}{3!} + \cdots .
\]
\end{definition}
The series converges for every $A$ and $t$. Differentiating term by term,
\[
   \frac{\d}{\d t}e^{At} = A + A^{2}t + \frac{A^{3}t^{2}}{2!}+\cdots
   = A\Bigl(I+At+\frac{(At)^{2}}{2!}+\cdots\Bigr) = Ae^{At},
\]
so $x(t)=e^{At}x_0$ satisfies $\dot x=Ax$ and $x(0)=x_0$: it is \emph{the}
solution. The properties follow the scalar ones, with one exception:
\begin{equation}\label{eq:expmprops}
   e^{A\cdot0}=I,\qquad e^{At_1}e^{At_2}=e^{A(t_1+t_2)},\qquad
   \bigl(e^{At}\bigr)^{-1}=e^{-At},\qquad
   e^{At}e^{Bt}\neq e^{(A+B)t}\ \text{unless } AB=BA .
\end{equation}
The third follows from the second with $t_2=-t_1$; in particular $e^{At}$ is always
invertible. The last fails because the series multiply out with $AB$ and $BA$
appearing in different orders. In MATLAB the matrix exponential is
\verb|expm|, not \verb|exp| (which exponentiates entry by entry): the position of
the second cart at $t=10$ from $x_0$ is \verb|[0 0 1 0]*expm(A*10)*x0|.

\subsection{The state response}\label{sec:stateresponse}
With an input, multiply $\dot x-Ax=Bu$ by $e^{-At}$. The left side becomes an exact
derivative, $e^{-At}\dot x-e^{-At}Ax=\frac{\d}{\d t}\bigl(e^{-At}x(t)\bigr)$, and
integrating from $0$ to $t$ gives $e^{-At}x(t)-x_0=\int_0^{t}e^{-A\tau}Bu(\tau)\,\d\tau$.
Multiplying by $e^{At}$:
\begin{equation}\label{eq:voc}
   \boxed{\;x(t) = e^{At}x_0 + \int_{0}^{t} e^{A(t-\tau)}\,B\,u(\tau)\,\d\tau ,\qquad
   y(t) = Ce^{At}x_0 + C\!\int_{0}^{t} e^{A(t-\tau)}Bu(\tau)\,\d\tau + Du(t).\;}
\end{equation}
This is the \emph{variation of constants} formula. Its first term is the
\emph{zero-input response}, what the initial state does on its own; the second is
the \emph{zero-state response}. The matrix $e^{A(t-t_0)}$ is the \emph{state
transition matrix}: it carries the state at time $t_0$ to the state at time $t$
when no input acts, $x(t)=e^{A(t-t_0)}x(t_0)$. The integral says the same about
the input: the small push $Bu(\tau)\,\d\tau$ delivered at time $\tau$ is carried
forward to time $t$ by $e^{A(t-\tau)}$, and all the pushes are added up.

Taking Laplace transforms of $\dot x=Ax+Bu$ instead gives
$X(s)=(sI-A)^{-1}x_0+(sI-A)^{-1}BU(s)$, which is \eqref{eq:voc} term by term
--- a convolution in time is a product in $s$. Comparing the two,
\begin{equation}\label{eq:laplaceexpm}
   \Lap\{e^{At}\} = (sI-A)^{-1}.
\end{equation}
The zero-state part of the output, $C(sI-A)^{-1}BU+DU$, is the transfer function of
\S\ref{sec:ss2tf} once more; the zero-input part is the free response that the
transfer function leaves out.

\subsection{Modes: solving by diagonalization}\label{sec:modes}
If $A=V\Lambda V^{-1}$ is diagonalizable, change to the coordinates
$z=V^{-1}x$ of \S\ref{sec:diag}. Then $\dot z=\Lambda z$ is $n$ separate scalar
equations $\dot z_i=\lambda_iz_i$, so $z_i(t)=e^{\lambda_it}z_i(0)$, and going back,
\begin{equation}\label{eq:modal}
   x(t) = Vz(t) = z_1(0)\,e^{\lambda_1t}\,v_1 + \dots + z_n(0)\,e^{\lambda_nt}\,v_n .
\end{equation}
Each term is a \emph{mode}: a fixed direction $v_i$, scaled by a scalar exponential.
The eigenvector decides the direction, the eigenvalue decides how fast that
component grows or decays, and the initial condition decides how much of each mode
is present (Figure~\ref{fig:modes}). The same formula, written without coordinates,
uses the \emph{left eigenvectors} $w_i\tp$ --- the rows of $V^{-1}$, which satisfy
$w_i\tp A=\lambda_iw_i\tp$ and $w_i\tp v_j=\delta_{ij}$:
\begin{equation}\label{eq:expmdyad}
   e^{At} = Ve^{\Lambda t}V^{-1} = \sum_{i=1}^{n} e^{\lambda_it}\,v_iw_i\tp ,
   \qquad z_i(0) = w_i\tp x_0 .
\end{equation}
So $e^{At}$ is a sum of rank-one matrices, one per mode.

\begin{keybox}
\textbf{Stability of $\dot x=Ax$.} $x(t)\to0$ for every initial state if and only
if every eigenvalue of $A$ has strictly negative real part.
\end{keybox}
For diagonalizable $A$ this is visible in \eqref{eq:modal}: a complex pair
$\sigma\pm j\omega$ contributes $e^{\sigma t}$ times a rotation at frequency
$\omega$, and every term dies out exactly when every $\operatorname{Re}\lambda_i<0$.
\S\ref{sec:jordanresp} shows that the statement survives without
diagonalizability. It is the state-space twin of the stable-plant condition of
Definition~\ref{def:stable}, and since the eigenvalues of $A$ are the poles of $P$
(barring cancellations), the two say the same thing.

\begin{figure}[htb]
\centering
\begin{tikzpicture}[>=Stealth,scale=1.35]
  \draw[->,gray] (-1.4,0) -- (2.6,0) node[right,lbl] {$x_1$};
  \draw[->,gray] (0,-0.9) -- (0,2.8) node[above,lbl] {$x_2$};
  \draw[dashed,gray] (-1.2,-0.48) -- (2.5,1.0);
  \draw[dashed,gray] (0.30,-0.86) -- (-0.92,2.63);
  \draw[->,very thick] (0,0) -- (1,0.4) node[below right,lbl] {$v_1$};
  \draw[->,very thick] (0,0) -- (-0.35,1) node[left,lbl] {$v_2$};
  \draw[->] (0,0) -- (2,0.8) node[right,lbl] {$z_1(0)v_1$};
  \draw[->] (0,0) -- (-0.525,1.5) node[left,lbl] {$z_2(0)v_2$};
  \draw[dotted] (2,0.8) -- (1.475,2.3) -- (-0.525,1.5);
  \fill (1.475,2.3) circle (1.4pt) node[above right,lbl] {$x(0)$};
  \draw[thick] plot[domain=0:9,samples=90,smooth,variable=\t]
    ({2*exp(-0.5*\t)-0.525*exp(-2*\t)},{0.8*exp(-0.5*\t)+1.5*exp(-2*\t)});
  \pgfmathsetmacro{\ax}{2*exp(-0.5*0.45)-0.525*exp(-2*0.45)}
  \pgfmathsetmacro{\ay}{0.8*exp(-0.5*0.45)+1.5*exp(-2*0.45)}
  \pgfmathsetmacro{\bx}{2*exp(-0.5*0.55)-0.525*exp(-2*0.55)}
  \pgfmathsetmacro{\by}{0.8*exp(-0.5*0.55)+1.5*exp(-2*0.55)}
  \draw[->,thick] (\ax,\ay) -- (\bx,\by);
  \fill (0,0) circle (1.2pt);
\end{tikzpicture}
\caption{A modal decomposition, drawn for eigenvalues $\lambda_1=-0.5$ and
$\lambda_2=-2$. The initial state splits along the eigenvectors (dotted
parallelogram); each component then shrinks along its own eigenvector at its own
rate. The fast mode $v_2$ dies first, so the trajectory bends and comes in to the
origin along the slow direction $v_1$.}
\label{fig:modes}
\end{figure}

\begin{example}[one matrix exponential, three ways]\label{ex:expm3}
Let $A=\begin{bmatrix}0&1\\-2&-3\end{bmatrix}$, with
$\det(\lambda I-A)=\lambda^{2}+3\lambda+2=(\lambda+1)(\lambda+2)$.

\emph{By diagonalization.} Eigenvectors $v_1=(1,-1)\tp$ for $\lambda_1=-1$ and
$v_2=(1,-2)\tp$ for $\lambda_2=-2$, so
\[
   V=\begin{bmatrix}1&1\\-1&-2\end{bmatrix},\quad
   V^{-1}=\begin{bmatrix}2&1\\-1&-1\end{bmatrix},\quad
   e^{At}=e^{-t}\begin{bmatrix}2&1\\-2&-1\end{bmatrix}
         +e^{-2t}\begin{bmatrix}-1&-1\\2&2\end{bmatrix},
\]
the two rank-one pieces $v_1w_1\tp$ and $v_2w_2\tp$ of \eqref{eq:expmdyad}.

\emph{By Laplace.} $(sI-A)^{-1}=\dfrac{1}{(s+1)(s+2)}\begin{bmatrix}s+3&1\\-2&s\end{bmatrix}$,
and the partial fractions of \S\ref{sec:pfe} entry by entry, e.g.\
$\frac{s+3}{(s+1)(s+2)}=\frac{2}{s+1}-\frac{1}{s+2}$, invert to the same matrix.

\emph{By Cayley--Hamilton} (\S\ref{sec:cayley}). Write $e^{At}=\alpha_0(t)I+\alpha_1(t)A$
and impose $e^{\lambda t}=\alpha_0+\alpha_1\lambda$ at both eigenvalues:
$e^{-t}=\alpha_0-\alpha_1$, $e^{-2t}=\alpha_0-2\alpha_1$. Hence
$\alpha_1=e^{-t}-e^{-2t}$, $\alpha_0=2e^{-t}-e^{-2t}$, and
$\alpha_0I+\alpha_1A$ is once more the matrix above.
\end{example}

\subsection{When there are too few eigenvectors}\label{sec:jordanresp}
For a matrix in Jordan form the decoupling is incomplete but still easy. With the
$3\times3$ $J$ of \S\ref{sec:jordan}, $\dot z=Jz$ reads
\[
   \dot z_1=\lambda_1z_1,\qquad \dot z_2=\lambda_2z_2+z_3,\qquad \dot z_3=\lambda_2z_3 .
\]
The first and last are solved as before. The middle one is driven by $z_3$, and by
\eqref{eq:voc} with $n=1$,
\[
   z_2(t) = e^{\lambda_2t}z_2(0)+\int_0^{t}e^{\lambda_2(t-\tau)}e^{\lambda_2\tau}z_3(0)\,\d\tau
          = e^{\lambda_2t}z_2(0) + t\,e^{\lambda_2t}z_3(0).
\]
Hence
\[
   x(t) = z_1(0)e^{\lambda_1t}v_1 + \bigl(z_2(0)e^{\lambda_2t}+z_3(0)\,t\,e^{\lambda_2t}\bigr)v_2
          + z_3(0)e^{\lambda_2t}v_3 ,
\]
and for a single block $J_2(\lambda)$, $e^{J_2t}=e^{\lambda t}\begin{bmatrix}1&t\\0&1\end{bmatrix}$.
A Jordan block brings a factor $t$ in front of the exponential --- the same $t$
that a repeated pole brought in \S\ref{sec:pfe}. Since $t^{k}e^{\lambda t}\to0$
whenever $\operatorname{Re}\lambda<0$, the stability criterion above holds for
every square matrix.

\subsection{The Cayley--Hamilton theorem}\label{sec:cayley}
\begin{theorem}[Cayley--Hamilton]
Every square matrix satisfies its own characteristic equation: if
$p(\lambda)=\det(\lambda I-A)=\lambda^{n}+a_{n-1}\lambda^{n-1}+\dots+a_0$, then
\[
   p(A) = A^{n}+a_{n-1}A^{n-1}+\dots+a_1A+a_0I = 0 .
\]
\end{theorem}
\begin{proof}[Sketch]
If $A=V\Lambda V^{-1}$, then $A^{k}=V\Lambda^{k}V^{-1}$, so
\[
   p(A)=V\,p(\Lambda)\,V^{-1},\qquad
   p(\Lambda)=\diag\bigl(p(\lambda_1),\dots,p(\lambda_n)\bigr)=0,
\]
because each $\lambda_i$ is a root. For a general $A$ the same argument runs block by block on
the Jordan form, where $p$ contains a factor $(\lambda-\lambda_i)^{k}$ large enough
to annihilate each $k\times k$ block.
\end{proof}
The consequence is that $A^{n}$ is a combination of $I,A,\dots,A^{n-1}$, hence so is
$A^{n+1}=A\cdot A^{n}$, and so on: \emph{every polynomial in $A$ reduces to one of
degree at most $n-1$}, and so does every convergent power series --- in particular
$e^{At}=\alpha_0(t)I+\alpha_1(t)A+\dots+\alpha_{n-1}(t)A^{n-1}$, as used in
Example~\ref{ex:expm3}. When the eigenvalues are distinct, the coefficients come
from the $n$ scalar equations $e^{\lambda_it}=\sum_k\alpha_k(t)\lambda_i^{k}$.

\begin{example}[reducing a matrix polynomial]
Let $A=\begin{bmatrix}2&1\\1&2\end{bmatrix}$, so $p(\lambda)=\lambda^{2}-4\lambda+3$ and
$A^{2}=4A-3I$. Then $A^{3}=4A^{2}-3A=13A-12I$ and $A^{4}=13A^{2}-12A=40A-39I$, so
\[
   A^{4}+3A^{3}+A^{2}+I = 83A-77I .
\]
By the same token $A^{10000}=aA+bI$ for two numbers $a,b$; evaluating the scalar
identity $\lambda^{10000}=a\lambda+b$ at the eigenvalues $1$ and $3$ gives
$a=\frac{3^{10000}-1}{2}$, $b=\frac{3-3^{10000}}{2}$.
\end{example}

\begin{remark}[four ways to compute $e^{At}$]
The series definition; the inverse Laplace transform of $(sI-A)^{-1}$
\eqref{eq:laplaceexpm}; the modal sum $\sum_ie^{\lambda_it}v_iw_i\tp$ when $A$ is
diagonalizable (with the Jordan form when it is not); and the Cayley--Hamilton
polynomial of degree $n-1$. They are four views of one object, and
Example~\ref{ex:expm3} checks three of them against each other.
\end{remark}

%%%%%%%%%%%%%%%%%%%%%%%%%%%%%%%%%%%%%%%%%%%%%%%%%%%%%%%%%%%%%%%%%%%%%%%%%%%%%%
\section{Positive Definiteness and the Singular Value Decomposition}\label{sec:svd}

\subsection{Quadratic forms and positive definite matrices}\label{sec:pd}
A \emph{quadratic form} is $x\tp Px$ for a square $P$. Only the symmetric part of
$P$ contributes, since $x\tp Px$ is a scalar and so equals its own transpose
$x\tp P\tp x$:
\[
   x\tp Px = x\tp\Bigl(\tfrac{P+P\tp}{2}\Bigr)x .
\]
It is therefore no loss to take $P=P\tp$.
\begin{definition}[positive definite]
A symmetric $P$ is \emph{positive definite}, written $P\succ0$, if $x\tp Px>0$ for
every $x\neq0$, and \emph{positive semidefinite}, $P\succeq0$, if $x\tp Px\ge0$
for every $x$.
\end{definition}
Positive definite matrices are the matrices that measure size: energies, costs and
squared distances are all quadratic forms of this kind.

\begin{proposition}[two tests]\label{prop:pdtests}
For $P=P\tp\in\R^{n\times n}$ the following are equivalent:
\begin{enumerate}[label=(\alph*),itemsep=2pt]
\item $P\succ0$;
\item every eigenvalue of $P$ is positive;
\item $P=L\tp L$ for some $L$ with full column rank.
\end{enumerate}
\end{proposition}
\begin{proof}
By Theorem~\ref{thm:spectral}, $P=V\Lambda V\tp$, so with $y=V\tp x$,
$x\tp Px=y\tp\Lambda y=\sum_i\lambda_iy_i^{2}$. As $x$ ranges over the nonzero
vectors so does $y$, and the sum is positive for all of them exactly when every
$\lambda_i>0$: (a)$\Leftrightarrow$(b). If $P=L\tp L$, then $x\tp Px=\|Lx\|^{2}\ge0$,
with equality only when $Lx=0$, i.e.\ only when $x=0$ if $\mathcal N(L)=\{0\}$:
(c)$\Rightarrow$(a). Conversely, (b) lets one take $L=\Lambda^{1/2}V\tp$.
\end{proof}
Semidefiniteness is the same with ``positive'' replaced by ``nonnegative'' and
the rank condition on $L$ dropped. Geometrically, for $P\succ0$ the level set
$x\tp Px=1$ is an ellipsoid whose axes point along the eigenvectors of $P$, with
semi-axis $1/\sqrt{\lambda_i}$ along $v_i$. For $P=\begin{bmatrix}2&1\\1&2\end{bmatrix}$
the eigenvalues are $1$ and $3$, so $P\succ0$, and $2x_1^{2}+2x_1x_2+2x_2^{2}=1$ is
an ellipse with its long axis along $(1,-1)$.

\subsection{The four fundamental subspaces}\label{sec:foursub}
Every $A\in\R^{m\times n}$ of rank $r$ carries four subspaces, two in each of the
spaces it connects:
\begin{center}\small
\renewcommand{\arraystretch}{1.3}
\begin{tabular}{@{}llll@{}}
\toprule
subspace & lives in & dimension & description\\
\midrule
$\mathcal R(A\tp)$ & $\R^{n}$ & $r$   & row space\\
$\mathcal N(A)$    & $\R^{n}$ & $n-r$ & null space\\
$\mathcal R(A)$    & $\R^{m}$ & $r$   & column space\\
$\mathcal N(A\tp)$ & $\R^{m}$ & $m-r$ & left null space\\
\bottomrule
\end{tabular}
\end{center}
The pairs are not merely of complementary dimension; they are perpendicular. If
$Ax=0$ and $y=A\tp z$, then $y\tp x=z\tp Ax=0$, so
\[
   \mathcal N(A)\perp\mathcal R(A\tp),\qquad\text{and likewise}\qquad
   \mathcal N(A\tp)\perp\mathcal R(A).
\]
Every $x\in\R^{n}$ thus splits uniquely as $x=x_r+x_n$ with $x_r$ in the row space
and $x_n$ in the null space, and $Ax=Ax_r$: the null-space part is lost, and $A$
maps the row space one-to-one onto the column space (Figure~\ref{fig:foursub}).

\begin{figure}[htb]
\centering
\begin{tikzpicture}[>=Stealth,scale=1.15]
  % ---- R^n ----
  \node[lbl] at (-1.7,1.9) {$\R^{n}$};
  \draw[thick] ($(0,0)+(215:1.3)$) -- ($(0,0)+(35:2.0)$);
  \draw[thick] ($(0,0)+(125:1.5)$) -- ($(0,0)+(305:1.4)$);
  \draw ($(35:0.18)$) -- ($(35:0.18)+(125:0.18)$) -- ($(125:0.18)$);
  \node[lbl,anchor=south west] at (35:2.0) {$\mathcal R(A\tp)$, dim $r$};
  \node[lbl,anchor=south east] at (125:1.5) {$\mathcal N(A)$, dim $n-r$};
  \fill (0,0) circle (1.2pt);
  \fill (35:1.35) circle (1.4pt) node[above left,lbl] {$x_r$};
  \fill (305:1.0) circle (1.4pt) node[left,lbl] {$x_n$};
  % ---- R^m ----
  \begin{scope}[shift={(6.2,0)}]
    \node[lbl] at (2.1,1.9) {$\R^{m}$};
    \draw[thick] ($(215:1.4)$) -- ($(35:2.0)$);
    \draw[thick] ($(125:1.3)$) -- ($(305:1.4)$);
    \draw ($(35:0.18)$) -- ($(35:0.18)+(125:0.18)$) -- ($(125:0.18)$);
    \node[lbl,anchor=south west] at (35:2.0) {$\mathcal R(A)$, dim $r$};
    \node[lbl,anchor=north west] at (305:1.4) {$\mathcal N(A\tp)$, dim $m-r$};
    \fill (0,0) circle (1.2pt) node[below left,lbl] {$0$};
    \fill (35:1.55) circle (1.4pt) node[above left,lbl] {$Ax_r$};
  \end{scope}
  \draw[->] ($(35:1.35)+(0.08,0.02)$) to[bend left=12]
        node[above,lbl] {$A$} ($(6.2,0)+(35:1.55)+(-0.08,0)$);
  \draw[->] ($(305:1.0)+(0.08,-0.02)$) to[bend right=14]
        node[below,lbl] {$A$} ($(6.2,0)+(-0.06,-0.04)$);
\end{tikzpicture}
\caption{The four fundamental subspaces of $A\in\R^{m\times n}$, each drawn as a
line. In $\R^{n}$ the row space and the null space are perpendicular, and so are the
column space and the left null space in $\R^{m}$. $A$ carries the row-space part
$x_r$ of a vector one-to-one onto the column space and sends the null-space part
$x_n$ to zero.}
\label{fig:foursub}
\end{figure}

One more fact links $A$ to the symmetric matrix $A\tp A$:
\begin{equation}\label{eq:nullata}
   \mathcal N(A\tp A) = \mathcal N(A),\qquad\text{hence}\qquad \rank(A\tp A)=\rank A .
\end{equation}
Clearly $Ax=0$ implies $A\tp Ax=0$. Conversely, if $A\tp Ax=0$ then
$0=x\tp A\tp Ax=\|Ax\|^{2}$, so $Ax=0$. In particular $A\tp A$ is invertible
exactly when $A$ has full column rank --- and then $A\tp A\succ0$, by
Proposition~\ref{prop:pdtests}(c) with $L=A$. In every case $A\tp A\succeq0$.

\subsection{The singular value decomposition}\label{sec:svdthm}
Eigenvectors only make sense for square matrices, and even then they need not be
orthogonal or even numerous enough. The singular value decomposition repairs all of
this at once, for every matrix.
\begin{theorem}[SVD]\label{thm:svd}
Every $A\in\R^{m\times n}$ of rank $r$ can be written
\[
   A = U\Sigma V\tp ,
\]
where $U\in\R^{m\times m}$ and $V\in\R^{n\times n}$ are orthogonal ($U\tp U=I$,
$V\tp V=I$) and $\Sigma\in\R^{m\times n}$ is zero except for the diagonal entries
$\sigma_1\ge\sigma_2\ge\dots\ge\sigma_r>0$. The $\sigma_i$ are the \emph{singular
values}, the columns $v_i$ of $V$ the right singular vectors, and the columns $u_i$
of $U$ the left singular vectors. Equivalently,
\begin{equation}\label{eq:svddyad}
   A = \sum_{i=1}^{r}\sigma_i\,u_iv_i\tp ,\qquad Av_i=\sigma_iu_i,\qquad A\tp u_i=\sigma_iv_i .
\end{equation}
\end{theorem}
\begin{proof}[Construction]
$A\tp A$ is symmetric and positive semidefinite, so by
Theorem~\ref{thm:spectral} it has orthonormal eigenvectors $v_1,\dots,v_n$ with
eigenvalues $\lambda_1\ge\dots\ge\lambda_n\ge0$, and by \eqref{eq:nullata}
exactly $r$ of them are positive. Put $\sigma_i=\sqrt{\lambda_i}$ and
$u_i=Av_i/\sigma_i$ for $i\le r$. These are orthonormal:
\[
   u_i\tp u_j=\frac{v_i\tp A\tp Av_j}{\sigma_i\sigma_j}
            =\frac{\lambda_j\,v_i\tp v_j}{\sigma_i\sigma_j}=\delta_{ij}.
\]
Complete $u_1,\dots,u_r$ to an orthonormal basis of $\R^{m}$. Then $AV=U\Sigma$,
column by column ($Av_i=0$ for $i>r$, since those $v_i$ lie in
$\mathcal N(A\tp A)=\mathcal N(A)$), and multiplying by $V\tp$ on the right gives
$A=U\Sigma V\tp$.
\end{proof}
So the singular values are
\[
   \sigma_i(A) = \sqrt{\lambda_i(A\tp A)} ,
\]
and the singular vectors are orthonormal bases for all four subspaces of
\S\ref{sec:foursub} at once: $v_1,\dots,v_r$ span the row space and
$v_{r+1},\dots,v_n$ the null space; $u_1,\dots,u_r$ span the column space and
$u_{r+1},\dots,u_m$ the left null space. The rank is the number of nonzero
singular values. The sum \eqref{eq:svddyad} is reading (iii) of \S\ref{sec:vspace},
$AB=\sum a_ib_i\tp$, with the factors chosen orthonormal and ordered by size.

\begin{example}[a $2\times2$ SVD]\label{ex:svd}
Let $A=\begin{bmatrix}3&0\\4&5\end{bmatrix}$. Then
$A\tp A=\begin{bmatrix}25&20\\20&25\end{bmatrix}$, with eigenvalues $45$ and $5$ and
eigenvectors $v_1=\tfrac{1}{\sqrt2}(1,1)\tp$, $v_2=\tfrac1{\sqrt2}(1,-1)\tp$. So
$\sigma_1=3\sqrt5\approx6.71$ and $\sigma_2=\sqrt5\approx2.24$, and
\[
   u_1=\frac{Av_1}{\sigma_1}=\frac{1}{\sqrt{10}}\begin{bmatrix}1\\3\end{bmatrix},\qquad
   u_2=\frac{Av_2}{\sigma_2}=\frac{1}{\sqrt{10}}\begin{bmatrix}3\\-1\end{bmatrix},
\]
which are indeed orthogonal. As a check,
$\sigma_1u_1v_1\tp+\sigma_2u_2v_2\tp=\tfrac32\begin{bmatrix}1&1\\3&3\end{bmatrix}
+\tfrac12\begin{bmatrix}3&-3\\-1&1\end{bmatrix}=A$, and
$\sigma_1\sigma_2=15=|\det A|$. The eigenvalues of $A$ are $3$ and $5$ --- not the
singular values. For a non-symmetric matrix the two are different objects; for a
symmetric one the singular values are the absolute values of the eigenvalues.
\end{example}

\subsection{What the singular values measure}\label{sec:svduse}
Read $A=U\Sigma V\tp$ from right to left: $V\tp$ rotates the input (it is
orthogonal, so it preserves lengths), $\Sigma$ stretches the $i$-th coordinate by
$\sigma_i$, and $U$ rotates the result. Hence $A$ maps the unit sphere onto an
ellipsoid whose semi-axes are $\sigma_iu_i$, and the input directions that end up
on those axes are the $v_i$ (Figure~\ref{fig:svd}).

\begin{figure}[htb]
\centering
\begin{tikzpicture}[>=Stealth]
  % ---- input: unit circle ----
  \begin{scope}[scale=1.15]
    \draw[->,gray] (-1.5,0) -- (1.6,0); \draw[->,gray] (0,-1.5) -- (0,1.6);
    \draw[thick] (0,0) circle (1);
    \draw[->,very thick] (0,0) -- (45:1) node[above right,lbl] {$v_1$};
    \draw[->,very thick] (0,0) -- (-45:1) node[below right,lbl] {$v_2$};
    \node[lbl] at (-1.15,1.35) {$\|x\|=1$};
  \end{scope}
  \draw[->,thick] (2.2,0.35) -- node[above,sig] {$A$} (3.4,0.35);
  % ---- output: ellipse A(circle) ----
  \begin{scope}[xshift=5.6cm,scale=0.33]
    \draw[->,gray] (-4.6,0) -- (4.8,0); \draw[->,gray] (0,-7.2) -- (0,7.6);
    \draw[thick] plot[domain=0:360,samples=120,smooth,variable=\a]
      ({3*cos(\a)},{4*cos(\a)+5*sin(\a)});
    \draw[->,very thick] (0,0) -- (2.121,6.364) node[right,lbl] {$\sigma_1u_1$};
    \draw[->,very thick] (0,0) -- (2.121,-0.707) node[right,lbl] {$\sigma_2u_2$};
  \end{scope}
\end{tikzpicture}
\caption{The SVD of Example~\ref{ex:svd} as a picture. The unit circle (left) is
carried by $A$ onto an ellipse (right, drawn at one third the scale) with semi-axes
$\sigma_1u_1$ and $\sigma_2u_2$, of lengths $3\sqrt5$ and $\sqrt5$. The input
directions $v_1$ and $v_2$ are the ones that land on the axes.}
\label{fig:svd}
\end{figure}

Several useful numbers follow directly.
\begin{itemize}[leftmargin=*,itemsep=2pt]
\item \textbf{Gain.} For every $x\neq0$,
$\sigma_{\min}\le\|Ax\|/\|x\|\le\sigma_1$ (with $\sigma_{\min}=\sigma_n$ if
$A$ has full column rank, and $0$ otherwise). The largest amplification is
$\sigma_1$, attained for $x$ along $v_1$; it is the \emph{induced $2$-norm}
$\|A\|_2=\sigma_1$. So singular values say how much a matrix can amplify, and in
which input direction --- where eigenvalues say how fast each mode evolves.
\item \textbf{Conditioning.} The ratio $\kappa(A)=\sigma_1/\sigma_r$ is the
\emph{condition number}. When it is large, some input directions are amplified far
more than others, and solving $Ax=b$ magnifies errors in $b$ by up to $\kappa$. A
matrix that is invertible on paper but has $\sigma_r$ tiny is, numerically, nearly
rank-deficient.
\item \textbf{Low-rank approximation.} Truncating \eqref{eq:svddyad} after $k$
terms, $A_k=\sum_{i\le k}\sigma_iu_iv_i\tp$, gives the closest matrix of rank $k$
to $A$, with error $\|A-A_k\|_2=\sigma_{k+1}$. The singular values say exactly how
much is lost by keeping only the strongest directions.
\end{itemize}

\begin{remark}[least squares]
When $Ax=b$ has no solution ($b\notin\mathcal R(A)$), the SVD still gives the best
one can do. The \emph{pseudo-inverse} $A^{+}=V\Sigma^{+}U\tp$, where $\Sigma^{+}$
inverts the nonzero singular values and transposes the shape, yields
$x=A^{+}b$, which minimizes $\|Ax-b\|$ and, among all minimizers, has the smallest
norm. If $A$ has full column rank, $A^{+}=(A\tp A)^{-1}A\tp$, which is invertible by
\eqref{eq:nullata}.
\end{remark}

\end{document}
